% Shape versus circuit: graphene geometric diodes and Tesla valves as
% terahertz rectifiers.
% Draft 2026-11-22 (eleventh draft: the lead lengths, the thin-gate ladders and the 1 um valve of Stage H8,
% results/HYDRO_REFEREE.md; tenth draft 2026-11-21: corrections after the review in paper/review_ninth_draft.md, sections 1 to 3;
% ninth draft 2026-11-21: the estimate of the drift limit and the results of Stage H7, results/HYDRO_CLOSEOUT.md;
% eighth draft 2026-11-20: text and figures corrected after the review in
% paper/review_seventh_draft.md; small-signal results of the 1 um funnel from the 20 nm mesh,
% results/thz2/fine, numbers as band ranges from tools/paper_numbers.py; seventh draft 2026-11-19:
% the fluid model with the momentum equation of the degenerate
% Dirac gas, Stage H6, results/HYDRO_MOMENTUM.md; sixth draft 2026-11-17, text corrections after
% the review in paper/review_fifth_draft.md; fifth draft 2026-11-16,
% reflecting contacts and device D with the field computed, Stage B10; fourth draft 2026-11-15, rough edges, Stage B9;
% third draft 2026-11-12, the short ballistic funnel with the field
% computed, Stage B8; second draft 2026-11-11 after the prior-work search and
% Stages B7 and H5). Every number is taken from results/*.md of the repository (commit
% to be tagged at submission); the figures are produced by
% paper/make_figures.py. Build: see paper/README.md.
\documentclass[aps,prapplied,twocolumn,superscriptaddress,floatfix,longbibliography]{revtex4-2}
\usepackage{graphicx}
\usepackage{amsmath,amssymb}
\usepackage{bm}
\usepackage{array}
\newcolumntype{L}[1]{>{\raggedright\arraybackslash}p{#1}}
\usepackage[colorlinks=true,allcolors=blue]{hyperref}

\newcommand{\vF}{v_{F}}
\newcommand{\EF}{E_{F}}
\newcommand{\kB}{k_{B}}
\newcommand{\RI}{R_{I}}
\newcommand{\Rant}{R_{\mathrm{ant}}}
\newcommand{\Rdc}{R_{\mathrm{dc}}}
\newcommand{\Pav}{P_{\mathrm{av}}}
\newcommand{\Idc}{I_{\mathrm{dc}}}
\newcommand{\taumr}{\tau_{\mathrm{mr}}}
\newcommand{\tauee}{\tau_{ee}}
\newcommand{\Rey}{\mathrm{Re}}
\newcommand{\Da}{\mathrm{Da}}
\newcommand{\um}{\,\mu\mathrm{m}}
\newcommand{\nm}{\,\mathrm{nm}}
\newcommand{\THz}{\,\mathrm{THz}}
\newcommand{\GHz}{\,\mathrm{GHz}}
\newcommand{\AW}{\,\mathrm{A/W}}
\newcommand{\pWrtHz}{\,\mathrm{pW/Hz^{1/2}}}
\newcommand{\percm}{\,\mathrm{cm^{-2}}}
\newcommand{\figpath}{figs}

\raggedbottom
\begin{document}

\title{Shape versus circuit: graphene geometric diodes and Tesla valves as terahertz rectifiers}

% AUTHOR BLOCK: affiliation and contact to be supplied by the author.
\author{Todd Chapman}
%\affiliation{}
%\email{}

\date{\today}

\begin{abstract}
Geometric diodes, asymmetric channels meant to rectify by their shape, and
electronic Tesla valves have both been proposed as junction-free rectifiers
for terahertz rectennas in graphene. We compute the rectification of one
set of funnel and Tesla-valve geometries in the ballistic, viscous and
plasmonic transport limits and separate what the shape contributes from
what the circuit contributes. The Tesla valve does not rectify usefully
in any case computed: its ballistic dc asymmetry in an imposed field
appears only in the clean specular limit and has the wrong sign, its
viscous diodicity is at most three per
cent at the drift velocities the graphene Fermi liquid allows, and at
terahertz frequencies the shape of the fluid valve contributes at most
$0.02\AW$ at either of two sizes; as a ballistic device it was not
computed at terahertz frequencies. In the
gated electron fluid, for which we use the momentum equation of the
degenerate Dirac gas, the funnel rectifies mainly through its
terminations: between $10\GHz$ and $1\THz$ a $1\um$ funnel with leads
of $12\um$ on a
$75\ \Omega$ antenna with the Dyakonov-Shur termination responds at
$0.05$ to $0.07\AW$ with the parameters of a fluid at room temperature
and at $0.14$ to $0.26\AW$ with those of one at $100$~K, which are at the
edge of the fluid regime and give more than $0.18\AW$ only if the
viscosity falls with frequency; by its shape alone it responds at
$0.10\AW$ or less. How the two compare at room temperature depends on
the leads: between $0.3$ and $1\THz$ the shape gives $2$ to $4\%$ of the
response with leads of $12\um$ and $4$ to $54\%$, typically two fifths,
with leads of $1\um$, shorter than the decay length of the plasmon.
The response of the funnel follows a square law up to a drift speed of three tenths of
the Fermi velocity, the limit to which we trust the model, and that speed
is reached at $26$ to $130$ microwatts. The ballistic funnel does rectify by its shape. In a uniform
imposed field its response depends on $fL/\vF$ alone and is halved near
$0.4$, which would limit the $6\um$ funnel with the best documented dc
asymmetry to tens of gigahertz; the static field that we compute for that
device is concentrated at its neck, however, and its terahertz response
with a computed field remains open. For a $0.55\um$ funnel we compute the
field in a particle simulation with a gate, metal contacts and
displacement current, without collisions between electrons. At $1\THz$
it rectifies $0.12$~A per absorbed watt,
five to seven times less than imposed fields give, and only if its
contacts absorb electrons arriving obliquely and its edges are specular
enough: half the response is lost at a specularity of $0.84$ if that is
the same at all angles, and at a head-on specularity of $0.5$ if grazing
electrons are reflected specularly, as a model of rough edges predicts.
Under a gate at a twentieth of the length of the device the gate
connection rectifies more than the shape at $1\THz$.
At the drift limit the fluid element converts $10^{-5}$ to
$1.2\times10^{-4}$ of the available power in the time domain; the square
law reaches $5\times10^{-4}$ at $100$~K only with a viscosity that falls
with frequency, which the time-domain solver does not have; the ballistic
funnel converts $2\times10^{-5}$ of the absorbed power at $0.1$~V and
$1.5\times10^{-4}$ at $0.2$~V. A matched
array of short funnels that meets those conditions would be a square-law
detector with a Johnson-limited noise-equivalent power of the order of
$100$~pW/Hz$^{1/2}$ at room temperature; it is not an energy harvester.
\end{abstract}

\maketitle

\section{Introduction}
\label{sec:intro}

A rectenna converts radiation to direct current by feeding an antenna into a
rectifier. Above a few hundred gigahertz the rectifier is the difficulty: the
$RC$ time of a junction diode limits Schottky rectifiers to the low terahertz
range, and the search for a junction-free, low-capacitance nonlinearity that
works to $1\THz$ and beyond is long standing. Graphene has attracted two
proposals of this kind, both of which rest on the shape of a channel rather
than on a junction.

The first is the geometric diode~\cite{moddel2012ultrahigh,zhu2013graphene,dragoman2013geometrically}: a channel
that narrows asymmetrically, typically a funnel, whose dimensions are below the
mean free path so that carriers are steered ballistically by the walls. With
planar geometry and negligible capacitance it was proposed explicitly for
terahertz and infrared rectennas, and the proposal came with a measurement:
a $0.5\um$ diode with a $75\nm$ neck in a bow-tie antenna gave a rectified
current at zero bias under $28\THz$ illumination~\cite{zhu2013graphene}.
Encapsulated devices reach reported dc asymmetry ratios of $1.6$ to
$2$~\cite{nguyen2021geometric,wang2024graphene}. Related four-terminal
ballistic rectifiers have been used for detection and imaging up to
$0.7\THz$~\cite{auton2016graphene,auton2017terahertz} and recently at
$3\THz$~\cite{shi2026ballistic}, and funnel-shaped rectifiers of chemical
vapor deposited graphene respond at zero bias from $0.49$ to
$0.68\THz$~\cite{hemmetter2026enhanced}; geometric diodes of silicon
nanowires have been measured with a flat response to
$40\GHz$~\cite{custer2020ratcheting}. On the side of theory, a
time-dependent particle simulation of a diode of about $100\nm$ found that
its rectification persists into the terahertz range and disappears above
about $5\THz$, a cut-off the authors place near the Drude damping
rate~\cite{stearns2020high}, and self-consistent simulations of ballistic
funnelling in graphene resonators find nonlinear response from below $10$ to
tens of terahertz~\cite{do2025nonlinear}. What these results leave open is
what sets the upper frequency of a given geometric diode, and whether the
long devices with wide leads, which show the best dc asymmetry, share the
speed of the short ones.

The second proposal is the electronic Tesla valve. Tesla's valvular
conduit~\cite{tesla1920valvular} rectifies fluid flow by
inertia~\cite{nguyen2021early}, and in the temperature window where
electron-electron collisions make the graphene electron liquid
viscous~\cite{bandurin2016negative,krishnakumar2017superballistic} a
Tesla-shaped channel has been reported to show a five per cent asymmetry of
its four-terminal resistance~\cite{geurs2020rectification}. A compressible
hydrodynamic simulation of a graphene Tesla valve gave a diodicity of
$1.07$~\cite{gungor2022modeling}, and in a GaAs electron gas a Tesla valve
has reached a diodicity of about $40$ above a threshold
current~\cite{sarypov2026electron}. A shape that rectifies an electron fluid
with no junction is an obvious rectenna candidate.

What the shape can do is constrained from the outset. For an elastic
two-terminal conductor the unitarity of the scattering matrix implies equal
transmission in the two directions at every energy, so the linear response
is symmetric whatever the shape. Rectification is a second-order effect,
and a consistent theory of it requires the response of the potential
landscape inside the device to the
bias~\cite{christen1996gauge,sanchez2004magnetic,song1999formalism,lofgren2004symmetry};
an asymmetric lowering of a barrier by the bias has indeed been proposed as
the origin of rectification in geometric diodes in place of ballistic
steering~\cite{bai2021asymmetric}. The useful question is therefore not
whether a shape rectifies but which nonlinearity it feeds, how strong that
nonlinearity is at the power an antenna delivers, and how fast it is. A
geometric diode in the ballistic regime feeds the bias-induced landscape; a
Tesla valve in the viscous regime feeds the convective term of the
Navier-Stokes equation; and at terahertz frequencies a gated graphene
channel is a plasmonic cavity whose convective and density nonlinearities
rectify by the mechanism of Dyakonov and
Shur~\cite{dyakonov1993shallow,dyakonov1996detection}, which is the basis of
most graphene terahertz detectors to
date~\cite{vicarelli2012graphene,bandurin2018resonant}. For field-effect
detectors of this kind the conversion efficiency is known to be quadratic
in the amplitude at small signal, to saturate when the amplitude reaches
the gate swing~\cite{gutin2012plasmonic,rudin2014theory,but2014nonlinear},
and to have an absolute maximum of the order of ten per cent, reached below
threshold~\cite{kachorovskii2013performance}; a silicon transistor rectenna
has delivered $13.6\%$ at $263\GHz$~\cite{khan2022cmos}. These mechanisms
have been studied separately, with different geometries and different
figures of merit. Whether a geometric shape, as opposed to the terminations
of the channel and its connection to the gate, carries the response has not
been tested with one set of devices in all of them.

This paper does that. We take one set of geometries
(Fig.~\ref{fig:devices}): the funnel of the best-documented geometric
diode~\cite{nguyen2021geometric}, the classical Tesla valve at the same
minimum feature size, a straight control, versions scaled to a $1\um$
channel that a $75\ \Omega$ antenna can drive, and four shorter funnels
down to $0.5\um$. We compute their rectification (i) in a semiclassical
ballistic model of the kind used in the geometric-diode
literature~\cite{zhu2013graphene,stearns2021simulation,truccolo2024comprehensive}, (ii) in
steady viscous electron hydrodynamics, without and with the gate, (iii) in
the gated compressible (plasmonic) electron fluid from $10\GHz$ to $3\THz$,
at small signal and at finite amplitude up to the gate overdrive, (iv)
in the ballistic model driven from $0.3\GHz$ to $30\THz$, and (v), for
the shortest funnel, in a particle simulation that computes the field
from the charge of the electrons, a gate and the contacts. Each solver is
validated against analytic limits and benchmarks before it is used; the
acceptance tests are listed in the Supplemental Material.

The answer separates cleanly into shape and circuit, and it differs
between the two devices. The Tesla valve does not rectify ballistically,
rectifies viscously only at Reynolds numbers the graphene Fermi liquid
cannot sustain, and contributes $10^{-5}$ to $10^{-3}\AW$ at terahertz
frequencies. In the electron fluid the funnel does no better as a shape:
with the gate included, a steady current sees a percent-level asymmetry
that is the same for every channel and belongs to the gate connection, and
the terahertz response of $0.1$ to $0.3\AW$ is the Dyakonov-Shur asymmetry
of the terminations, with a conversion efficiency below one per cent at
every operating point for which the fluid description holds. The ballistic
funnel is different. Its shape does rectify at terahertz frequencies, on
conditions. The first is its length: the response is a function of the
reduced frequency $fL/\vF$ and is gone when the flight through the device
takes longer than about two periods. The $6\um$ device of the dc
experiments fails that condition at $1\THz$; a $0.55\um$ funnel meets it.
The second is the field. With an imposed potential, the simplest
assumption of simulations of geometric diodes, the short funnel appears to give several
times the current per absorbed watt of the plasmonic element; with the
potential computed it gives a sixth of that, about what the plasmonic
element gives at room temperature. The third is the edges. A rough edge
reflects oblique electrons more specularly than head-on ones, and the
walls of the taper are met obliquely; with such edges the response is
halved at a head-on specularity of $0.5$ and gone at $0.2$, and etched
edges have been measured from near zero to
$0.8$~\cite{taychatanapat2013electrically,masubuchi2012boundary,barnard2017absorptive,berdyugin2020minibands}.
The fourth is the contacts, and their resistance is a poor guide to it.
A contact that reflects three quarters of the electrons, whatever their
angle, costs the short funnel far less than the same resistance in
series, while a contact of equal resistance that admits only electrons
near normal incidence removes the rectification. And when the gate is closer than
the device is long, the gate connection again rectifies more than the
shape. As a detector the short funnel with specular edges and contacts
that absorb at all angles is where the
plasmonic element and the best measured graphene detectors are. As a harvester
neither element exists: no mechanism studied here converts as much as one
per cent of the incident power.

Section~\ref{sec:models} describes the devices and the models,
Sec.~\ref{sec:dc} the dc results, Sec.~\ref{sec:plasmonic} the plasmonic
element from small signal to saturation, Sec.~\ref{sec:ballistic} the
ballistic element at terahertz frequencies, its scaling with length and
the short funnel with a computed field, and
Sec.~\ref{sec:discussion} compares the regimes, places the numbers against
measured detectors and states the limitations.

\begin{figure*}
\includegraphics[width=\textwidth]{\figpath/fig1_devices.pdf}
\caption{Devices and frequency scales. (a) The funnel geometric diode,
device D of Ref.~\cite{nguyen2021geometric}: $2\um$ leads, a 30 degree taper
to a $175\nm$ neck and an abrupt shoulder; forward is the direction into the
taper. (b) The classical four-loop Tesla valve at the same $175\nm$ channel
width. (c) The gated funnel scaled to a $1\um$ neck with $4\um$ contacts, with
the terminations used for the plasmonic response: the antenna, a Thevenin
source of resistance $\Rant$, drives the source contact and the drain is tied
to the gate. (d) Characteristic frequencies of the mechanisms studied, for the
parameters of Table~\ref{tab:params}; the shaded band is the range of
interest for a terahertz rectenna. The shorter funnels of
Sec.~\ref{sec:length} are drawn in Fig.~\ref{fig:length}(a).}
\label{fig:devices}
\end{figure*}

\section{Devices and models}
\label{sec:models}

\subsection{Geometries}

All devices are two-terminal planar channels with the source contact on the
left and the drain on the right. The funnel [Fig.~\ref{fig:devices}(a)] has
the dimensions of device D of Nguyen \textit{et al.}~\cite{nguyen2021geometric}:
leads $2\um$ wide, a taper at 30 degrees to the axis down to a neck of
$175\nm$, and an abrupt shoulder on the drain side, $6\um$ from end to end,
with the contacts on the two ends. Its measured two-terminal asymmetry,
the ratio of the voltages at $\pm80\ \mu$A, is $1.66$; the device has a
resistance of $12$ to $20$~k$\Omega$ and carries that current at about
$1.5$~V. The classical Tesla valve [Fig.~\ref{fig:devices}(b)] is
built from channel centre lines: a zig-zag main channel of width $175\nm$
inclined at $\pm 25$ degrees, and at every vertex a loop that continues the
arriving segment straight on, turns through a half circle of radius $175\nm$
and rejoins the main channel against the stream. Flow from left to right
follows the zig-zag; flow from right to left is split into the four loops. We
also keep three earlier variants in which the loops are attached to one side
of a straight channel, because they illustrate that the conclusion does not
depend on the details of the valve, and a straight $175\nm$ channel between
the same leads as a control. For the terahertz calculations each shape is
also scaled so that its narrowest width is $1\um$, with $4\um$ contacts,
which brings the device resistance from kilohms to the hundred-ohm range of
an antenna [Fig.~\ref{fig:devices}(c)]; the scaled funnel has straight
leads of $12\um$ on each side of its taper and is $27\um$ long. Four further funnels, from a
half-scale copy of device D down to a $0.5\um$ device with the proportions
of the diode measured at $28\THz$~\cite{zhu2013graphene}, are used for the
length scaling of Sec.~\ref{sec:length} and listed in
Table~\ref{tab:short}.

\subsection{Ballistic model}
\label{sec:model-ballistic}

The ballistic model is the semiclassical trajectory (Monte
Carlo~\cite{jacoboni1983monte}) method of the geometric-diode
literature~\cite{zhu2013graphene,stearns2021simulation,truccolo2024comprehensive}.
Electrons with the Dirac dispersion $E=\hbar\vF|\bm{k}|$,
$\vF=10^{6}$~m/s, are injected at the contacts with the flux-weighted cosine
law, fly along classical trajectories, reflect from the walls specularly with
probability $p$ and diffusely otherwise, and suffer isotropic
momentum-relaxing collisions after exponentially distributed free paths of
mean $\lambda$. The current is the Landauer expression
\begin{equation}
I=\frac{g e}{h}\int dE\,N(E)\,T(E;V)\,\bigl[f_{S}(E)-f_{D}(E)\bigr],
\label{eq:landauer}
\end{equation}
with $g=4$, $N(E)=EW/\pi\hbar\vF$ the number of modes of the injecting contact
of width $W$, $T$ the transmission probability estimated from the
trajectories, and $f_{S,D}$ the Fermi functions of the contacts at
electrochemical potentials $\EF$ and $\EF-eV$. Energies are sampled from the
bias window $N|f_{S}-f_{D}|$, so the estimator has zero variance at zero bias.

With no potential inside the device, flux conservation gives
$N_{S}T_{SD}=N_{D}T_{DS}$ at every energy and Eq.~\eqref{eq:landauer} is odd
in $V$ for any shape; the trajectories cannot rectify. Rectification enters
through the potential energy $U(\bm{r};V)$ that the bias creates inside the
channel~\cite{christen1996gauge}. We use the simplest assumption of the published
simulations, a uniform field along the channel, $U=-eVx/L$ over the device
length $L$, and integrate the curved flights in chords with a controlled
sagitta. Self-consistent electrostatics is known to roughly halve the
asymmetry of such diodes relative to the uniform-field value and to
reverse it for shallow tapers~\cite{truccolo2024comprehensive}; we return
to the field distribution in Sec.~\ref{sec:ballistic}. The figure of merit is the asymmetry $A=I(+V)/|I(-V)|$; writing
$I=G_{1}V+G_{2}V^{2}$, $A\simeq 1+2(G_{2}/G_{1})V$, and a cosine drive of
amplitude $V_{a}$ rectifies $\Idc=G_{2}V_{a}^{2}/2$ if the response is
instantaneous. Carriers are electrons at a chemical potential
$\EF=0.138$~eV and $300$~K. This is the Fermi energy of a density of
$1.4\times10^{12}\percm$, which we estimate from the gate voltage and
lever arm reported for that device~\cite{nguyen2021geometric}, an upper
bound because the gate voltage of its neutrality point is not reported;
at $300$~K the electron density at this chemical potential, which is the
density of the simulated gas, is $1.56\times10^{12}\percm$. Holes give
the inverse asymmetry. The model has no length of its own: with the
bias, $\EF$ and $\kB T$ fixed, the transmission depends only on the shape
and on $\lambda/L$, and under a drive of frequency $f$ only on those and on
$fL/\vF$, while currents and conductances scale with the width. This is
used, and tested, in Sec.~\ref{sec:length}.

The solver was validated on eleven tests before use: unit transmission of a
specular channel, the free-molecular transmission factors of a diffuse
channel~\cite{richley1964numerical}, the Sharvin
conductance~\cite{sharvin1965possible}, the ballistic-to-diffusive crossover, reciprocity with and
without the field, unit asymmetry of a mirror-symmetric constriction, and
convergence of the chord integration. Classical trajectories require
features well above the Fermi wavelength, $30\nm$ here. Collisions
between electrons are not in this model; their mean free path at $300$~K
is of the order of $100\nm$ (Table~\ref{tab:params}), shorter than
every device studied, and we return to them in Secs.~\ref{sec:short} and
\ref{sec:discussion}.

\subsection{Ballistic model with a computed field}
\label{sec:model-pic}

For the $0.55\um$ funnel G3 of Table~\ref{tab:short}, and for device D
at dc, the potential is computed instead of imposed, in a particle simulation of the kind applied to geometric diodes in
Refs.~\cite{stearns2020high,stearns2021simulation}. The graphene sheet
lies in a uniform dielectric of permittivity $3.8$ at a distance $d$ above
a grounded gate plane; $d$ is $30\nm$, the gate of the plasmonic model
below, or $300\nm$, an oxide back gate that leaves a device of this size
nearly ungated. The contacts are metal strips in the plane of the sheet,
$200\nm$ long and as wide as the leads, held at potentials $V_{S}(t)$ and
$V_{D}(t)$. All electrons of the device, about $3200$, are represented by
$4$ to $8\times10^{6}$ sample particles that move together, with the
trajectories, wall reflection, scattering and equilibrium injection at the
contacts of the model above. At every time step of $2.5$~fs their charge
is assigned to a grid of $5\nm$ spacing and the potential energy
\begin{equation}
U(\bm{r})=\frac{e^{2}}{4\pi\epsilon}\sum_{\bm{r}'}\delta N(\bm{r}')
\left[\frac{1}{|\bm{r}-\bm{r}'|}-\frac{1}{\sqrt{|\bm{r}-\bm{r}'|^{2}+4d^{2}}}\right]
\label{eq:kernel}
\end{equation}
is evaluated by fast Fourier transform; $\delta N$ is the electron number
at a grid node in excess of the uniform background that makes the undriven
sheet neutral, the second term is its image in the gate, and the sum
includes the charge on the contacts, which is solved for at every step so
that the strips are equipotentials. The particles are then advanced in
the gradient of the interpolated potential. This is the electrostatics of
a thin sheet in three dimensions; it goes over to the local gate
capacitance of the plasmonic model for small $d$ and to the ungated sheet
for large $d$, and it contains the plasma waves of the collisionless gas.

The currents at the source, the drain and the gate include the
displacement current, so the admittance $Y_{1}$ at the drive frequency and
the absorbed power $P_{\mathrm{abs}}=\langle V_{S}I_{S}+V_{D}I_{D}\rangle$
are results of the calculation. With the gate present the terminations
matter, as in the plasmonic model, and two are used: a drive that leaves
only the shape, $V_{S}=-V_{D}=(V_{a}/2)\cos\omega t$, under which a
mirror-symmetric device gives no rectified current by symmetry, and the
drain tied to the gate, $V_{S}=V_{a}\cos\omega t$ and $V_{D}=0$, under
which a straight channel already rectifies.

The walls reflect either with one specularity $p$ for all angles of
incidence, as in the model above, or as a rough edge. For the rough edge
an electron of wave number $k$ that meets the wall at an angle $\theta$
from its normal is reflected specularly with probability
\begin{equation}
p(\theta)=\exp\left[-(2hk\cos\theta)^{2}\right],
\label{eq:rough}
\end{equation}
where $h$ is the rms height of the
roughness~\cite{soffer1967statistical,shao2017probing}: the more obliquely
an electron arrives, the more specular the edge. The other electrons
leave at an angle $\theta'$ drawn from the density
$[1-p(\theta')]\cos\theta'$. This choice makes the wall reciprocal, so
that a gas in equilibrium stays uniform and isotropic; re-emission by the
cosine law alone does not, once $p$ depends on the angle (Supplemental
Material). We label a rough edge by its specularity for a head-on hit at
the Fermi wave vector $k_{F}=\EF/\hbar\vF=0.21\nm^{-1}$,
$p_{0}=\exp[-(2hk_{F})^{2}]$, which is $0.5$ for $h=2.0\nm$.

The contacts described so far are ideal: the line on which the sheet
meets a strip absorbs every electron that reaches it and injects the
equilibrium flux of the strip. Two contacts that reflect are also used.
In the first an arriving electron is absorbed with a probability $T$,
the same for all angles, and reflected as by a mirror otherwise; the
injected flux is reduced by the factor $T$, which keeps an undriven gas
in equilibrium. A wide, clean strip between two such contacts has
$T/(2-T)$ of its ballistic conductance. A transfer-length measurement
would assign each contact $(2-T)/T$ times the $48\ \Omega\um$ that an
ideal contact shows at this density, half of the ballistic resistance of
the sheet: $81$, $145$ and $338\ \Omega\um$ for $T=0.75$, $0.5$ and
$0.25$. In the second the probability depends on the angle $\theta$
between the electron and the normal to the contact line as it does for a
smooth $p$-$n$ junction of half width
$d_{j}$~\cite{cheianov2006selective},
\begin{equation}
T(\theta)=\exp\left(-\pi k d_{j}\sin^{2}\theta\right),
\label{eq:collimate}
\end{equation}
so that electrons near normal incidence pass and the others are
reflected, and the injected flux is reduced by the same function. Such a
junction stands in front of a contact whose metal dopes the graphene
beside it to the polarity opposite to that of the channel. We use
$d_{j}=7$ and $38\nm$, for which the strip conducts as between contacts
of $T=0.5$ and $0.25$. For device D the same model is run with ideal
contacts, the gate at $300\nm$, a grid of $7.5\nm$ and $5\times10^{7}$
particles for its $1.6\times10^{5}$ electrons.

Tests T22 to T32 of the Supplemental Material were passed before use,
two of them in a form that had to be relaxed and is described there.
They include the
capacitance of the contact strips, the static screening of a weak
potential, the speed of the gated plasma wave, which is $0.98$ to $0.99$
of the collisionless value of $3.90\,\vF$ under the $30\nm$ gate, the
results of the imposed-field model for the same device when the
interaction is switched off, a linear conductance that the interaction
leaves unchanged within $1\%$, the reflection law and the detailed
balance of the rough wall, and the absorption law, the strip conductance
and the equilibrium of the reflecting contacts. The sample particles collide with each
other through the grid, which relaxes the Fermi distribution towards a
classical one; in the open device a particle is replaced by a contact
within about a picosecond, and the rectified current at the operating
point of Sec.~\ref{sec:pic} does not change between $4$ and
$16\times10^{6}$ particles, while it is too large by up to $40\%$ with
$10^{6}$.

\subsection{Viscous model}
\label{sec:model-viscous}

Where electron-electron collisions are the fastest scattering process the
electron liquid obeys the Navier-Stokes
equation~\cite{gurzhi1968hydrodynamic,bandurin2016negative}. The dc model is
steady incompressible flow with momentum relaxation,
\begin{equation}
\Rey\,(\bm{u}\cdot\nabla)\bm{u}=-\nabla p+\nabla^{2}\bm{u}-\Da\,\bm{u},
\qquad \nabla\cdot\bm{u}=0,
\label{eq:ns}
\end{equation}
in units of the lead width $W$ and the mean inlet speed, with
\begin{equation}
\Rey=\frac{I}{n e \nu},\qquad \Da=\frac{W^{2}}{\nu\taumr},\qquad
S=\frac{\ell_{s}}{W}
\end{equation}
the Reynolds number, the damping (Gurzhi) number and the slip length
$\ell_{s}$ of the Navier wall condition~\cite{kiselev2019boundary} in units
of $W$. In two dimensions the Reynolds number depends on the current and not
on the width. The device is current biased with a developed inlet profile,
and the figure of merit is the diodicity $D=\Pi_{r}/\Pi_{f}$, the ratio of
the pressure (voltage) drops for flow from drain to source and from source
to drain. Equations~\eqref{eq:ns} are solved with Taylor-Hood
elements~\cite{taylor1973numerical} on triangular
meshes~\cite{geuzaine2009gmsh} with Newton iteration and a sparse direct
solver~\cite{guennebaud2010eigen}; the wall condition
$\bm{u}\cdot\bm{n}=0$ is imposed with length-weighted
normals~\cite{engelman1982implementation}, which makes the discrete wall flux
vanish exactly and removes a spurious diodicity of $10^{-5}$ that polygonal
walls otherwise produce. Validation covers Poiseuille and Brinkman flow with
slip, Stokes reciprocity ($D-1<10^{-10}$ at $\Rey=0$), the Kovasznay
solution~\cite{kovasznay1948laminar}, the backward-facing
step~\cite{armaly1983experimental,erturk2008numerical}, and the
diodicity of the valve itself as a classical fluid device. The discretisation
floor of $D-1$ is $3\times10^{-5}$.

\subsection{Plasmonic model}
\label{sec:model-plasmonic}

At terahertz frequencies the fluid is compressible: a gate at distance $d$
makes the local density follow the local electrochemical potential, and
density waves propagate as plasmons. With the gradual-channel approximation
and degenerate Dirac carriers the electrochemical potential is
$\Phi(n)=e^{2}n/C_{g}+\hbar\vF\sqrt{\pi n}$, the hydrodynamic mass is
$m^{*}(n)=\hbar\sqrt{\pi n}/\vF$~\cite{tomadin2013theory,svintsov2012hydrodynamic}, and because both
depend on the density only the force per unit mass is the gradient of a
single function $h(n)$ with $h'=\Phi'/m^{*}$. The model is
\begin{equation}
\partial_{t}n+\nabla\cdot(n\bm{u})=0,\label{eq:cont}
\end{equation}
\begin{multline}
\partial_{t}\bm{u}+(\bm{u}\cdot\nabla)\bm{u}
-\tfrac12\bm{u}\,(\nabla\cdot\bm{u})-\tfrac14\nabla u^{2}
+\frac{\bm{u}}{\taumr}\\
=-\nabla h(n)+\nu\nabla^{2}\bm{u},\label{eq:mom}
\end{multline}
with the plasma speed $s(n)$ given by
\begin{equation}
s^{2}(n)=n\,h'(n)=\frac{e^{2}n}{C_{g}m^{*}(n)}+\frac{\vF^{2}}{2}.
\label{eq:s2}
\end{equation}
The second term is the quantum capacitance~\cite{fang2007carrier}; its share $q=\vF^{2}/2s_{0}^{2}$
of $s_{0}^{2}=s^{2}(n_{0})$ is $0.04$ at $n_{0}=10^{12}\percm$ under $30\nm$ of
hBN and grows to $0.5$ at the thin-gate, low-density corner of
Sec.~\ref{sec:operating}. Replacing $h$ by $s_{0}^{2}n/n_{0}$ and dropping
the third and fourth terms on the left of Eq.~\eqref{eq:mom} gives the
shallow-water model of Dyakonov and Shur~\cite{dyakonov1993shallow},
which we call the classical gas; Eqs.~\eqref{eq:cont} and \eqref{eq:mom}
as written are the Dirac gas. In the small-signal solver the viscosity
has the viscoelastic form $\nu/(1+i\omega\tauee)$, which matters above
$1/\tauee$; the time-domain solver has the constant $\nu$.

The third and fourth terms on the left of Eq.~\eqref{eq:mom} are what
distinguishes the Dirac gas from an ordinary fluid in the inertia, and
they are of the same order as the convective term. For carriers of energy
$\hbar\vF|\bm{k}|$ that drift at a velocity $\bm{u}$, the momentum density
at fixed density is $n m^{*}\bm{u}$ and the momentum flux is
$P_{0}(n)+\tfrac34 n m^{*}u^{2}$ along the drift and
$P_{0}(n)-\tfrac14 n m^{*}u^{2}$ across it, where
$P_{0}=\hbar\vF\sqrt{\pi}\,n^{3/2}/3$ is the pressure of the gas at rest
and an ordinary fluid has
the factors $1$ and $0$; with Eq.~\eqref{eq:cont} this gives
Eq.~\eqref{eq:mom} to second order in the amplitude about a state of
uniform density at rest, which is the order to which the small-signal
response is computed. In one dimension the three
inertial terms cancel, so that the degenerate Dirac gas has no convective
term at
this order~\cite{svintsov2013hydrodynamic}. Two other forms are in the
literature. The form with
$(\bm{u}\cdot\nabla)\bm{u}$ alone, which is that of a classical electron
gas with the mass $m^{*}(n)$ and which Ref.~\cite{tomadin2013theory}
introduces as a crude approximation, gives a uniform channel between
$1.6$ and $3$ times the rectified voltage of Eq.~\eqref{eq:mom}
(Sec.~\ref{sec:small}); we refer to it as the ordinary convective term.
The photoresponse derived in Ref.~\cite{tomadin2013theory} rests on an
equation that keeps the third term of Eq.~\eqref{eq:mom} and not the
fourth, which leaves half a convective term in one dimension and gives
the uniform channel $1.3$ to $1.9$ times the voltage of
Eq.~\eqref{eq:mom} between $0.1$ and $1\THz$
(Appendix~\ref{app:perturbation}). We do not follow it: the general
equations of that reference reduce to Eq.~\eqref{eq:mom} once the pressure
is evaluated at fixed density in the frame of the lattice, and the speeds
$u_{0}/2\pm s_{0}$ of waves on a drifting
gas~\cite{svintsov2013hydrodynamic} require the cancellation. For an
incompressible flow the third term vanishes and the fourth is a gradient
that is absorbed in the pressure, so the viscous model of
Eq.~\eqref{eq:ns} holds for the Dirac liquid as it stands.
Equation~\eqref{eq:mom} is the expansion of the Euler equation of an
ideal relativistic fluid with $\vF$ in place of the speed of
light~\cite{landau1987fluid,tomadin2013theory}. For steady states and for
finite amplitudes we use that equation in full
[Eq.~\eqref{eq:exact} of Appendix~\ref{app:perturbation}], which holds for
any drift speed below $\vF$ and has no solution beyond it. Two choices
enter at the same order as these terms and are varied in
Sec.~\ref{sec:small}: a contact fixes the density of the fluid next to
it, as in
Refs.~\cite{dyakonov1993shallow,tomadin2013theory,svintsov2013hydrodynamic},
and the gas is taken at zero temperature.

The antenna is a Thevenin source of amplitude $V_{a}$ and resistance $\Rant$
connected to lumped ports: each contact is an equipotential whose voltage and
current satisfy Kirchhoff's law with the external circuit. Two terminations
are used throughout. In the \emph{symmetric} one the antenna drives the two
contacts in antiphase through $\Rant/2$ each and the gate sits at the
midpoint; the circuit is then mirror symmetric and any dc response is due to
the shape of the channel. In the \emph{Dyakonov-Shur} (DS) termination the
antenna drives the source through $\Rant$ and the drain is tied to the gate;
the terminations are then asymmetric and a uniform channel already
rectifies~\cite{dyakonov1996detection}. The responsivity is the
short-circuit dc current per available power,
\begin{equation}
\RI=\frac{I_{\mathrm{dc}}}{\Pav},\qquad \Pav=\frac{V_{a}^{2}}{8\Rant},
\end{equation}
and the Johnson-noise-limited noise-equivalent power is
$\mathrm{NEP}=\sqrt{4\kB T/\Rdc}/\RI$ with $\Rdc$ the dc resistance of the
device.

Two solvers are used. For small signals the fields are expanded to second
order about rest: the first order is a complex time-harmonic plasmon problem,
and the time average of the second order is a Stokes-Brinkman problem for the
dc flow, forced by products of first-order fields, whose port current or
voltage is the rectified signal (Appendix~\ref{app:perturbation}). For finite amplitudes
Eqs.~\eqref{eq:cont} and \eqref{eq:exact} are integrated in time with the
second-order backward difference formula and Newton iteration, with a bias
tee that separates the dc load from the antenna. Both use Taylor-Hood
elements on the meshes of the viscous model. The steady states of
Eqs.~\eqref{eq:cont} and \eqref{eq:exact} under a dc current, used in
Sec.~\ref{sec:gated}, are obtained with the same discretisation by Newton
iteration. Every finite-amplitude and steady solution carries its largest
drift speed, and results above $0.3\,\vF$ are marked: there the ideal gas
is still defined, but heating and the scattering by optical phonons,
which the model does not contain, are not negligible in a real sample.
We call $0.3\,\vF$ the drift limit. It is not a derived bound, and every
power quoted at the limit, and every square-law efficiency quoted there,
scales as its square. Three estimates place it (Supplemental Material).
First, the edge of a Fermi sea that drifts at a speed $u$ is spread in
energy by $2(u/\vF)\EF$ to lowest order in $u/\vF$, and the emission of an optical phonon
opens where that spread equals the phonon energy: at $0.41\,\vF$ for the
$100$~meV modes of hexagonal boron nitride~\cite{tielrooij2018out} at
$10^{12}$~cm$^{-2}$, and at higher speeds for lower densities and for the
phonons of graphene. Second, the absorbed power heats the electrons. With
the cooling time of $2.5$~ps measured in encapsulated graphene at room
temperature~\cite{tielrooij2018out}, the electrons of the reference
device at $300$~K are $5\%$ hotter than the lattice on the average over
the device when the largest drift is $0.3\,\vF$, and up to $30\%$
hotter near the source contact, where the drift is then largest
(Sec.~\ref{sec:large}). Below about $200$~K the measured
cooling is slower and no longer
exponential~\cite{tielrooij2018out}; with $10$~ps the average at
$100$~K is $20$ to $28\%$ above the lattice. Third, the drift velocity
measured in encapsulated graphene at room temperature saturates at
$0.6\,\vF$ at $8\times10^{11}$~cm$^{-2}$~\cite{yamoah2017high}, twice
the limit. At $0.3\,\vF$ the effects that the model omits are thus of
the order of tens of per cent. The limit is not conservative at
$100$~K, and these estimates do not fix it to better than between $0.2$
and $0.4\,\vF$.
The validation ladder has twenty-seven tests: the static limit against the viscous solver, the plasmon
dispersion, the Dyakonov-Shur closed form and its counterpart for the
Dirac gas [Fig.~\ref{fig:small}(a)],
reciprocity of the symmetric termination, mesh and time-step convergence,
agreement of the two solvers at small amplitude, invariance under
co-scaling of the physical parameters at fixed nondimensional groups,
the speeds of waves on a drifting fluid, which are $u_{0}/2\pm s_{0}$ for
the Dirac gas~\cite{svintsov2013hydrodynamic} against $u_{0}\pm s_{0}$ for
an ordinary one, a radial flow that separates the two new terms of
Eq.~\eqref{eq:mom}, and,
for the steady states, a quadrature for the gated channel and the
incompressible limit against the viscous solver.

\begin{table}
\caption{Parameter sets. The ballistic model follows the experiment of
Ref.~\cite{nguyen2021geometric}; the two plasmonic sets are representative of
hBN-encapsulated graphene at the two temperatures and are labelled by them;
they are assumptions, not fits. The $100$~K viscosity is of the order
measured~\cite{bandurin2016negative,krishnakumar2017superballistic}; the
$300$~K value is the kinetic estimate $\vF^{2}\tauee/4=0.025$~m$^{2}$/s,
rounded. The steady viscous results are given for the range $0.02$ to
$0.1$~m$^{2}$/s. The electron-electron mean free path $\vF\tauee$ is
$1\um$ at $100$~K and $0.1\um$ at $300$~K. $W$ is the contact width.}
\label{tab:params}
\begin{ruledtabular}
\begin{tabular}{ll}
Quantity & Value \\
\hline
\multicolumn{2}{l}{\textit{Ballistic}} \\
$\EF$, $T$ & $0.138$ eV, 300 K \\
$n$ at 300 K; at $T=0$ & $1.56$; $1.4\times10^{12}\percm$ \\
specularity $p$ & 1, 0.9, 0.8, 0.7, 0.5, 0 \\
mean free path $\lambda$ & $\infty$, 700, 200, 100 nm \\
\multicolumn{2}{l}{\textit{Viscous}} \\
current at $\Rey=1$ & $160\ \mu$A ($\nu=0.1$ m$^2$/s), $32\ \mu$A (0.02) \\
$\Da$, $S$ & 0 to 128; 0 to $\infty$ \\
\multicolumn{2}{l}{\textit{Plasmonic}} \\
$n_{0}$, $d$, $\epsilon_{r}$ & $10^{12}\percm$, 30 nm, 3.8 \\
$s_{0}$, gate overdrive & $3.6\,\vF$, 1.43 V \\
``100 K'': $\taumr$, $\tauee$, $\nu$ & 5 ps, 1 ps, 0.1 m$^2$/s \\
``300 K'': $\taumr$, $\tauee$, $\nu$ & 1 ps, 0.1 ps, 0.03 m$^2$/s \\
$W/s_{0}\taumr$ ($W=4\um$) & 0.22 (100 K), 1.12 (300 K) \\
antenna $\Rant$ & $75\ \Omega$ \\
\end{tabular}
\end{ruledtabular}
\end{table}

\section{Rectification at dc}
\label{sec:dc}

\subsection{Ballistic: the funnel rectifies, the Tesla valve does not act as a valve}

Figure~\ref{fig:ballistic-dc} shows the asymmetry at $0.1$~V and the
resistance of all six devices for three edge specularities and three mean
free paths. In the clean specular limit the funnel gives $A=1.62\pm0.02$
($1.25$ at $0.05$~V); at the mean free path estimated for these
devices~\cite{nguyen2021geometric}, $700\nm$, the same model gives $1.09$,
and $1.06$ with fully diffuse edges. The assumed field matters as much as
the scattering. If the same bias is taken to drop over the taper and neck
alone, with the leads cut to $250\nm$ (device G2 of
Table~\ref{tab:short}), the asymmetry at $\lambda=700\nm$ is $1.55$ and
the clean value $1.91$; if it drops over $500\nm$ centred on the neck of
the full device it is $0.96$, the reverse sign. The experiment has its
contacts at the ends of the $6\um$ shape and is a two-terminal
measurement, so none of these profiles is excluded by its geometry. The
choice of field profile, $\lambda$ and $p$ moves $A-1$ by a factor of
several and can reverse it.

Computing the field removes that choice. In the model of
Sec.~\ref{sec:model-pic}, with $\lambda=700\nm$, specular edges, ideal
contacts and the gate at $300\nm$, device D has
\begin{equation}
A=1.0174\pm0.0008
\label{eq:deviceD}
\end{equation}
at $0.1$~V [Fig.~\ref{fig:ballistic-dc}(a)], and $1.0172\pm0.0014$ with
half the number of particles; its resistance is $0.93$~k$\Omega$, with a
difference of $2.5\%$ between the two particle numbers that we have not
traced. The potential is that of neither assumption
(Fig.~\ref{fig:deviceD}). Of the bias, $55\%$ falls on the $500\nm$
centred on the neck, where a uniform field puts $8\%$; the two leads take
$9\%$ and $15\%$, the taper up to that region $15\%$, and $3\%$ falls at
each contact plane. The asymmetry lies between those of the two imposed
profiles, as the potential does, and $A-1$ is a fifth of the
uniform-field value at the same mean free path. Self-consistent
electrostatics was known to roughly halve the asymmetry of such
diodes~\cite{truccolo2024comprehensive}; for this shape it does more.

The comparison with the measurement has to be made at equal bias. The
asymmetry of $1.66$ reported for this device~\cite{nguyen2021geometric} is
the ratio of the voltages at $\pm80\ \mu$A, a current that the device
carries at about $1.5$~V. That is fifteen times the bias of the
calculations and about ten times $\EF/e$, so the nearness of the clean
uniform-field value of $1.62$ at $0.1$~V to it has no significance. At
$0.1$~V the published curve gives an asymmetry of $1.03$ to $1.04$, which
we read from the figure with an uncertainty of $0.02$
[Fig.~\ref{fig:ballistic-dc}(a)]. The computed $1.017$ is consistent with
it, but the comparison does not test the model: the measured device has a
resistance of $12$ to
$20$~k$\Omega$ against $0.93$~k$\Omega$, so most of its
resistance lies in what the model does not contain, and an
asymmetry this close to one would follow from many models. What the calculations
establish is the mechanism and the scale: electrons accelerated toward
the neck by the bias arrive with more kinetic energy and more transverse
modes than those decelerated toward it, and how much that rectifies
depends on where the bias drops.

No Tesla geometry acts as a valve in this model. All 54 points of the three
side-loop variants lie between $0.93$ and $1.08$ with standard errors of
$0.01$ to $0.05$, and sixteen of the eighteen points of the classical valve
are within two standard errors of one. The two exceptions are in the clean
specular limit, where the classical valve gives $A=0.93\pm0.02$ at $0.05$~V
and $0.89\pm0.01$ at $0.1$~V. That is as far from one as the funnel at
$\lambda=700\nm$, but of the wrong sign, the blocking direction being
the easier one, and it is gone as soon as edges are partly diffuse or bulk
scattering is present. The straight control gives $1.00\pm0.01$ and
$0.99\pm0.01$ in the clean specular limit; over all eighteen conditions
its scatter shows that the errors quoted for these runs are too small by
a factor of about $1.3$, which leaves the two exceptions of the valve
standing. These
errors are ten or more times those of the computed-field result of
Eq.~\eqref{eq:deviceD}, so an asymmetry of a per cent or two in a Tesla
valve is not excluded by them, and no Tesla valve was run with a computed
field or under an ac drive. Meanwhile
the loops cost resistance in both directions, $1.0$ to $22$~k$\Omega$ against
$0.5$ to $2.2$~k$\Omega$ for the funnel [Fig.~\ref{fig:ballistic-dc}(b)]. The
result follows from the structure of the problem rather than from the
parameters: at fixed potential landscape the transmission is reciprocal for
any shape, and the Tesla loops do nothing to the way the bias landscape acts
on the mode count at a constriction, which is what the funnel exploits. It agrees with the experiment, which
reports that the rectification of the valve disappears outside the viscous
regime~\cite{geurs2020rectification}.

\begin{figure*}
\includegraphics[width=\textwidth]{\figpath/fig2_ballistic_dc.pdf}
\caption{Ballistic transport at dc. (a) Asymmetry $A=I(+V)/|I(-V)|$ at
$0.1$~V for the funnel, four Tesla-valve geometries and the straight control,
for edge specularities $p$ and mean free paths $\lambda$; error bars are one
standard error from 64 batches of 400\,000 trajectories in total. The
black diamond is the funnel with the field computed,
Eq.~\eqref{eq:deviceD}; the grey bar is the asymmetry measured for the
funnel at $0.1$~V, read from the curve of
Ref.~\cite{nguyen2021geometric}. (b) Forward resistance of the same
devices in the imposed uniform field.}
\label{fig:ballistic-dc}
\end{figure*}

\begin{figure}
\includegraphics[width=\columnwidth]{\figpath/fig13_deviceD.pdf}
\caption{Device D with the field computed: static potential energy at a
bias of $0.1$~V in units of the bias, for $\lambda=700\nm$, specular
edges, ideal contacts and the gate at $300\nm$. (a) Map, with contours
every $0.1$ drawn after smoothing over $45\nm$; the bars are the contact
planes. (b) On the axis, against the two imposed profiles of the text.
The profile at the reversed bias differs from the one shown by less
than $0.04$, and by $0.001$ away from the neck.}
\label{fig:deviceD}
\end{figure}

\subsection{Viscous: the Tesla valve works, beyond the reach of the Fermi liquid}

In the electron fluid the picture reverses in principle and not in practice.
Figure~\ref{fig:viscous}(a) shows the diodicity of the $175\nm$ classical
valve and of the funnel against Reynolds number without damping. The valve is
a genuine Tesla valve: with slip walls $D-1$ is $0.15$ to $0.23$ at $\Rey=10$
and $1.1$ to $1.5$ at $\Rey=30$, growing close to $\Rey^{2}$, the signature of
an inertial loss present in one direction only. The fields
[Fig.~\ref{fig:viscous}(c) and \ref{fig:viscous}(d)] show its origin: in the
forward direction the loops are stagnant, in the reverse direction each
carries a stream that re-enters the main channel against the flow. As a
check on the geometry, the same valve treated as a classical no-slip fluid
device gives $D=1.05$ at $\Rey=50$ and $1.23$ at $\Rey=100$;
three-dimensional simulations of multistage valves give $1.02$ to $1.03$
and $1.10$ to $1.15$ at the same Reynolds
numbers~\cite{liosis2025diodicity}, a duct of finite depth adding wall
friction that the two-dimensional flow lacks. The diodicity of the valve is
converged with respect to the mesh to about $10\%$ (Supplemental
Material), which is immaterial for what follows. The funnel is a diode of the
opposite sign, the forward direction being the harder one: its asymmetry is
that of the exits, a jet from the abrupt shoulder against a diffuser along
the taper, and it is the stronger of the two below $\Rey\approx3$
($|D-1|=1.2\times10^{-2}$ at $\Rey=3$ with slip) but saturates near $0.15$.
This is the convective (Bernoulli) nonlinearity proposed for wedge-shaped
channels as a signature of electron
hydrodynamics~\cite{hui2021beyond} and recently measured in Venturi-shaped
GaAs devices~\cite{monari2026electron}; with equal lead widths on the two
sides, as here, it is quadratic in the current.

Three facts decide the matter. First, the experiment of
Ref.~\cite{geurs2020rectification} used currents of at most a few
microamperes ($6.3\ \mu$A for the device that showed the effect), for
which we estimate $\Rey\approx0.04$ to $0.2$ at $10^{12}\percm$ and
$\nu=0.1$ to $0.02$~m$^{2}$/s; the authors quote no Reynolds number. For
$\Rey$ up to $0.3$ every Tesla geometry, wall condition and damping gives
$|D-1|<3\times10^{-5}$ (the funnel, $1.5\times10^{-4}$), three orders of
magnitude below the reported five per cent, which was observed for holes in a single-loop valve as a ratio of
four-terminal resistances. Steady incompressible
electron hydrodynamics does not reproduce that observation; whatever
produces it is not the inertial Tesla mechanism. Second, the Reynolds number is bounded.
The drift speed in a channel of width $w$ is $u=\Rey\,\nu/w$, and keeping it
below $0.3\,\vF$, the drift limit that we assume for all the fluid models
(Sec.~\ref{sec:model-plasmonic}), caps
$\Rey$ at $0.5$ to $2.6$ in a $175\nm$ channel and at $3$ to $15$ in a
$1\um$ channel for $\nu=0.1$ to $0.02$~m$^{2}$/s [shaded in
Fig.~\ref{fig:viscous}]. The bound is not new: the Reynolds number of the
graphene electron liquid scales as the drift speed over the Fermi velocity
times the size of the flow over the electron-electron mean free path, and
values of order ten have been judged the most that can be
expected~\cite{muller2009graphene,lucas2018hydrodynamics}. What is added
here is what a Tesla valve delivers inside it. Third, the low viscosity needed for a large
Reynolds number comes with weak slip relative to the width and strong
relative damping. Figure~\ref{fig:viscous}(b) shows the valve scaled to
$1\um$: each factor of four in $\Da$ costs a factor two to three, and a
slip length of $0.4\um$ instead of perfect slip a factor twenty. A
Reynolds number of ten needs a viscosity of $0.02$~m$^{2}$/s or less,
which belongs to room temperature, and with the momentum-relaxation time
of that set $\Da$ is then $50$ in a $1\um$ channel: the curves of
Fig.~\ref{fig:viscous}(b) at smaller damping are not reachable at that
Reynolds number. Inside the
admissible window the diodicity of the $1\um$ valve is $10^{-3}$ to
$3\times10^{-2}$ over these dampings and slip lengths, at a
resistance of $0.3$ to $1.6$~k$\Omega$; the $175\nm$ valve costs $10$ to
$86$~k$\Omega$ for less.

These statements are for graphene, for steady laminar flow and below the
drift limit, and an experiment in another material shows what lies outside
them. In a GaAs electron gas at $4$~K
the forward and reverse characteristics of a Tesla valve coincide up to
about $300\ \mu$A and then separate, the reverse resistance rising to a
diodicity of about $40$, which the authors attribute to the onset of
turbulence at a Reynolds number of at least
$7$~\cite{sarypov2026electron}. Below the threshold this is the
reciprocity found here at low Reynolds number. Above it, the steady
solution of Eq.~\eqref{eq:ns} gives $D-1$ of order $0.1$ at $\Rey=7$ to
$10$ with slip walls, two orders of magnitude short: the large diodicity
is not the steady inertial loss computed in this
section, and unsteady, heated flow is outside Eq.~\eqref{eq:ns}.

\begin{figure*}
\includegraphics[width=\textwidth]{\figpath/fig3_viscous_dc.pdf}
\caption{Steady viscous flow. (a) Diodicity $|D-1|$ of the $175\nm$ classical
Tesla valve and of the funnel without damping, for three wall conditions;
filled symbols: reverse direction harder ($D>1$), open symbols: forward
harder. The bar marks the diodicity reported in
Ref.~\cite{geurs2020rectification} at our estimate of its Reynolds number; grey: below the discretisation floor;
orange: drift speed above $0.3\,\vF$. (b) The valve scaled to a $1\um$
channel, with momentum-relaxation damping $\Da$ and finite slip length; the
two shaded ranges are the drift-speed limits for $\nu=0.1$ and
$0.02$~m$^{2}$/s. (c), (d) Speed in the $175\nm$ valve at $\Rey=30$ with
perfect slip for forward and reverse flow.}
\label{fig:viscous}
\end{figure*}

\subsection{With the gate: a circuit term and no shape term}
\label{sec:gated}

Equation~\eqref{eq:ns} has no gate. In a gated device the density follows
the local potential, so a bias opens the channel at one end and closes it at
the other, and a second nonlinearity appears that owes nothing to inertia.
We computed the steady states of Eqs.~\eqref{eq:cont} and \eqref{eq:exact} for
both funnels, both Tesla valves, a straight channel and a mirror-symmetric
constriction, with a current source at one contact and the other contact
tied to the gate, for the two parameter sets of Table~\ref{tab:params} and
biases from $1.4$ to $430$~mV ($10^{-3}$ to $0.3$ of the gate overdrive
$U_{0}=e n_{0}/C_{g}=1.43$~V).

A two-terminal device with a gate can be reversed in two ways, and they
measure different things. Reversing the polarity of the source, with the
gate connection where it was, changes the sign of the bias relative to the
gate. Exchanging the contacts, so that the gate is tied to the other end
and the same current is sent through the shape in the other direction,
leaves the bias relative to the gate as it was and reverses only the flow
through the geometry. The first operation tests the circuit, the second the
shape. In a gradual channel of width $W(x)$, without inertia and viscosity,
the steady equations integrate to (Appendix~\ref{app:gated})
\begin{equation}
G(n_{s})-G(1)=\Gamma I\int_{0}^{L}\frac{dx}{W(x)},\qquad
G'(n)=n\,h'(n),
\label{eq:1d}
\end{equation}
in the units of Appendix~\ref{app:perturbation}, with $I$ the current and
$n_{s}$ the density at the driven contact. The shape enters through the
integral, which is the linear resistance, and through nothing else:
exchanging the contacts gives no asymmetry at all, and reversing the
polarity gives the same asymmetry for every shape,
\begin{equation}
D_{\mathrm{pol}}-1=\frac{V}{2V_{u}}=0.48\,\frac{V}{U_{0}}
\label{eq:dpol}
\end{equation}
to first order in the bias $V$ for the Dirac gas ($V/U_{0}$ for the
classical gas), the harder polarity being the one that depletes the channel. Here
$V_{u}=m^{*}s_{0}^{2}/e$ is the voltage unit of the gated fluid, $1.49$~V
at the reference point.

The two-dimensional solutions bear this out
[Fig.~\ref{fig:gated}(a)]. Under reversal of the polarity all six devices
lie on one curve, $1.45$ to $1.53\%$ at $43$~mV, within $5\%$ of the closed form
wherever the drift speed is below $0.3\,\vF$ ($8\%$ for the $1\um$ Tesla
valve at its largest bias inside the limit); the terms of relative order
$u^{2}/\vF^{2}$ in Eq.~\eqref{eq:exact}, which Eq.~\eqref{eq:1d} omits, and
the mesh contribute to the difference in comparable amounts (Supplemental
Material). Under exchange of the
contacts the asymmetry is at most $3\times10^{-4}$ at any point of the main set of runs
inside the drift limit. Where it is above the discretisation
floor it is the
inertial term of the previous subsection, little changed by the gate
[Fig.~\ref{fig:gated}(b)]: at nine of the ten
points at which the
incompressible diodicity exceeds $3\times10^{-5}$ the gated solution agrees with
it in magnitude within a factor two, and the exception is the $1\um$ Tesla valve at the edge of the
limit, where the gated value is ten times smaller. Beyond the limit the
comparison loses its meaning: the terms of Eq.~\eqref{eq:exact} beyond
second order grow, and a bias that would require a drift above $\vF$ has no
steady state at all ($12$ of the $112$ points computed). At low bias the exchange asymmetry on the meshes of the main set of runs is
linear in the bias, which Eq.~\eqref{eq:1d} forbids. It is numerical: for
the funnel its coefficient falls by a factor of sixteen as the wall mesh
is refined from $28$ to $10\nm$, while the quadratic, inertial coefficient
changes by a sixth to a quarter [Fig.~\ref{fig:gated}(c) and Supplemental
Material]. We therefore report the low-bias
shape term of a gated channel as a bound, $1.4\times10^{-5}$ at $14$~mV for
the funnel and $4.5\times10^{-7}$ for the Tesla valve.

Three things follow. A gated channel of any shape, a straight strip
included, shows an asymmetry of a per cent at some tens of millivolts when
its polarity is reversed, and a measurement or a simulation that reverses
the polarity will find it in a Tesla valve without the valve having
anything to do with it. The compressible simulation of
Ref.~\cite{gungor2022modeling}, which gives a diodicity of $1.07$ for a
graphene Tesla valve with a $10\nm$ channel at $50$~mV, is not reproduced
here as an effect of the shape: at $43$~mV the exchange asymmetry of our
valve is $10^{-5}$ on the mesh of the main runs and $10^{-6}$ on the finest
one, and the polarity term of Eq.~\eqref{eq:dpol} is
$1.7\%$ at $50$~mV for any channel under our gate. Second, the circuit term does not explain the measurement of
Ref.~\cite{geurs2020rectification}, which is a ratio of four-terminal
resistances and is reported to vanish outside the viscous regime. Third, the
result depends on the monolayer: the plasma speed under a $30\nm$ gate is
$3.6\,\vF$, so the drift limit keeps the flow below a tenth of the wave
speed. In bilayer graphene, where the flow can be driven to the wave
speed, a nozzle does show a polarity-dependent resistance of hydrodynamic
origin~\cite{geurs2025supersonic}, which is outside the present model.

The steady-flow asymmetries are also slow. The inertial one is established over the viscous
and convective times of the channel, $10$ to $20$~ps at $1\um$, so the
picture of this section applies below some tens of gigahertz. Above that
the response of a gated channel is carried by plasma waves, and the question
has to be asked of the time-dependent form of the gated model.

\begin{figure*}
\includegraphics[width=\textwidth]{\figpath/fig9_gated_dc.pdf}
\caption{Steady flow with the gate, $n_{0}=10^{12}\percm$ under $30\nm$ of
hBN. (a) Asymmetry of the voltage against bias for the $300$~K parameters
and slip walls, under reversal of the polarity (filled symbols) and under
exchange of the contacts (open symbols and crosses); the dashed line
follows from Eq.~\eqref{eq:1d}. The symmetric constriction under exchange
of the contacts measures the asymmetry of its unstructured mesh. (b)
Exchange asymmetry with the gate against the diodicity of the
incompressible model of Eq.~\eqref{eq:ns} at the same Reynolds number,
damping, wall condition and mesh, for both parameter sets; filled: slip
walls, open: no slip; the band is a factor two. (c) Refinement of the wall
mesh for the $175\nm$ funnel and Tesla valve at $300$~K, at biases of
$0.01\,U_{0}$ (filled symbols) and $0.03\,U_{0}$ (open symbols). Faint symbols in
(a) and (b): drift speed above $0.3\,\vF$.}
\label{fig:gated}
\end{figure*}


\section{The plasmonic element at terahertz frequencies}
\label{sec:plasmonic}

\subsection{Small signal: the terminations rectify, the shape tunes}
\label{sec:small}

The solver is first checked where an analytic answer exists. For a uniform
channel driven by an ideal source with an open drain, the classical gas
reproduces the closed form of Dyakonov and
Shur~\cite{dyakonov1996detection} to better than $0.1\%$ from $10\GHz$ to
$1.2\THz$ and to $0.6\%$ at $3\THz$, through the plasmon resonances at
$100$~K [Fig.~\ref{fig:small}(a)]. The Dirac gas of Eq.~\eqref{eq:mom}
without viscosity lies on its own closed form, Eq.~\eqref{eq:channel}, to
the same accuracy, and that form is exactly half the result of Dyakonov
and Shur at every frequency (Appendix~\ref{app:perturbation}). With the
viscosity of the $100$~K set the Dirac gas gives $0.43$ to $0.51$ of the
classical response from $10\GHz$ to $1\THz$.

Figure~\ref{fig:small}(b) shows the $1\um$ funnel, whose dc resistance
$\Rdc$ is $172\ \Omega$ at $100$~K and $582\ \Omega$ at $300$~K. With the
DS termination it absorbs $87$ to $100\%$ of the available power of a
$75\ \Omega$ source above $0.1\THz$ at either temperature, and its
responsivity is $0.14$ to $0.26\AW$ at $100$~K ($0.20$ to $0.26\AW$
between $0.3$ and $1\THz$) and $0.053$ to $0.066\AW$ at $300$~K, without
resonant structure from $10\GHz$ to $1\THz$. Above $1\THz$ the $100$~K
response is not converged with the mesh (Supplemental Material) and is
not shown; at $300$~K it falls slowly, to $0.044\AW$ at $3\THz$. With the
symmetric termination, which isolates the shape, the response at $100$~K
is $6\times10^{-4}$ to $0.02\AW$ below $100\GHz$ and rises to $0.06$ to
$0.10\AW$ between $0.2$ and $1\THz$ as the taper becomes comparable to the
plasmon wavelength. At $300$~K, where friction and viscosity damp the
plasmon, it is of order $10^{-4}\AW$ or less below $100\GHz$, which the
mesh does not resolve, and $0.5$ to $2.3$~mA/W between $0.2$ and $1\THz$.
Between $0.3$ and $1\THz$ the shape alone thus gives $0.27$ to $0.38$ of
the DS current at $100$~K, about a third, and $1.6$ to $4.2\%$ of it at
$300$~K. These are ratios of currents; in efficiency a third is a ninth.

The split at $300$~K is decided by the leads. The $1\um$ funnel has
straight leads of $12\um$ on each side, and a plasmon travels $3.6\um$ in
a momentum-relaxation time at $300$~K: with the symmetric drive the wave
launched at each contact has all but died out where the taper begins. We
repeated the sweeps with leads of $6$, $3$ and $1\um$ and nothing else
changed (Table~\ref{tab:leads}). The DS response rises by a factor of
$1.5$ to $3$ as the leads and their resistance are removed. The
shape-only response rises by a factor of $3$ to $90$, and its share of
the DS current between $0.3$ and $1\THz$ goes from $1.6$ to $4.2\%$ to
between $4$ and $54\%$, with a median over the six frequencies of the
sweep of $39\%$. With leads of $3$ and $1\um$ the shape-only response
changes sign inside the band, so that the share has deep minima; the
device is then $2.4$ and $1.3$ decay lengths long and resonates even at
$300$~K. The $175\nm$ funnel behaves in the same way: with leads of
$1\um$ it gives $-0.05\AW$ by its shape against $0.10\AW$ with the DS
termination at $1\THz$ and $300$~K. At $300$~K, then, the shape
contributes a few per cent in a device several decay lengths long and
something comparable to the terminations in one of about one decay
length.

The control is a plain channel with the width of the leads and the full
length of the funnel, on the same antenna. With the symmetric termination
it must give nothing, and it gives less than $3\times10^{-6}\AW$ up to
$1\THz$ at either temperature. With the DS termination at $300$~K and
leads of $12\um$ it gives $0.062$ to $0.066\AW$ between $0.3$ and
$1\THz$, and the funnel gives $0.83$ to $0.85$ of that at every frequency
from $0.1$ to $1\THz$: at room temperature the constriction takes away
from what the terminations give. The same holds with leads of $6$ and
$3\um$, where the funnel gives $0.73$ to $0.91$ of the channel of its own
length; with leads of $1\um$ the two have different resonances and the
funnel gives $0.8$ to $1.8$ of the channel. At no lead length does the
DS-terminated funnel do systematically better than a strip. At $100$~K
the channel with $12\um$ leads gives $0.11$ to $0.19\AW$ between $0.1$
and $1\THz$ and the funnel $1.1$ to $2.4$ times as much. At that
temperature the decay length is $18\um$ and every one of these devices is
a resonator, so that shortening the leads moves all the curves and tests
nothing cleanly: the share of a third is that of the $12\um$ leads, and
with shorter ones it lies between $0.02$ and $5$ from one frequency to
the next. With leads of $1\um$ the plain channel then gives $-0.05$ to
$-2.0\AW$ between $0.1$ and $1\THz$, larger than the funnel's in magnitude
at eight of ten frequencies and opposite in sign at eight (Supplemental
Material). We have not checked these $100$~K sweeps against the mesh or
examined the change of sign, and we take no number from them. The sign of
the DS response of the reference devices follows the termination:
exchanging the roles of source and drain reverses it. That the shape of
a channel modifies the response set by its boundary conditions, rather
than replacing it, is in line with what is known for field-effect
detectors: a non-uniform channel changes the responsivity by about ten
per cent~\cite{zhang2021terafet}, and a Corbino geometry, which
concentrates the field at the inner contact, raises it
substantially~\cite{khavronin2020singularity}.

\begin{table*}
\caption{The $1\um$ funnel at $300$~K against the length of its leads,
between $0.3$ and $1\THz$ (six frequencies) on a $75\ \Omega$ antenna.
$L/\ell$: length of the device over the $3.6\um$ that a plasmon travels
in a momentum-relaxation time. Share: magnitude of the shape-only
response over the DS response, range and median. Last column: DS
response of the funnel over that of a plain channel of the same length.}
\label{tab:leads}
\begin{ruledtabular}
\begin{tabular}{ccccccc}
Leads ($\mu$m) & $L/\ell$ & $\Rdc$ ($\Omega$) & DS (A/W) & Shape (mA/W) & Share (\%) & Funnel/channel \\
\hline
12 & 7.5 & 582 & 0.053--0.055 & 0.9--2.3 & 1.6--4.2 (2.4) & 0.83--0.85 \\
6 & 4.1 & 364 & 0.086--0.092 & 0.6--15 & 0.7--17 (12) & 0.75--0.83 \\
3 & 2.4 & 255 & 0.118--0.139 & 2.9--43 & 2.4--31 (10) & 0.75--0.91 \\
1 & 1.3 & 179 & 0.097--0.176 & 4.8--82 & 3.8--54 (39) & 0.79--1.8 \\
\end{tabular}
\end{ruledtabular}
\end{table*}

The $175\nm$ funnel [Fig.~\ref{fig:small}(c)] has a dc resistance of
$1.9$~k$\Omega$ and absorbs about half or more of the available power
above $0.1\THz$. Its shape does more than that of the wide device. At
$100$~K the shape-only response is $0.02$ to $0.19\AW$ in magnitude
between $0.3$ and $1\THz$, with changes of sign, against $0.06$ to
$0.30\AW$ with the DS termination. At $1.87\THz$ the taper passes through
a series resonance, the horn is matched to the antenna, and the
shape-only responsivity reaches $0.21\AW$ at $100$~K, half of the
$0.43\AW$ of the DS termination at the same frequency. A control with
$1\um$ leads has the resonance at the same frequency, so it belongs to
the taper. In a plasmonic design the funnel shape earns its keep in this
way, as a matched horn at low temperature. At $300$~K it gives
$0.013\AW$ against $0.066\AW$ for the terminations at the same frequency.

The Tesla valve [Fig.~\ref{fig:small}(d)] contributes little. The
shape-only responsivity of the $175\nm$ valve is below
$8\times10^{-4}\AW$ below $300\GHz$, where values under about
$5\times10^{-4}\AW$ are not resolved by the mesh, and reaches
$1.8\times10^{-2}\AW$ only at a $100$~K resonance near $0.7\THz$. At
$1\THz$ it is $1.0\times10^{-2}\AW$ at $100$~K and $4.5\times10^{-4}\AW$
at $300$~K, about twenty and thirty times below the $175\nm$ funnel. Its
dc resistance of $6$ to $11$~k$\Omega$ is two orders of magnitude from
the antenna, and what the DS termination adds, $0.004$ to $0.09\AW$, is
the terminations and not the loops. A valve with a channel of $1\um$,
every length scaled by $1000/175$, is nearer the antenna, at
$2.9$~k$\Omega$ at $300$~K and $0.9$~k$\Omega$ at $100$~K, and its shape
does no more (Supplemental Material). Between $0.1$ and $1\THz$, with
leads of $12\um$ or of $1\um$, its shape-only response is below
$9\times10^{-4}\AW$ at $300$~K and at most $0.021\AW$ at $100$~K. That is
$4$ to $1000$ times below the shape-only response of the $1\um$ funnel
with the same leads at $38$ of the $40$ frequencies and temperatures
compared; at the other two the funnel's response passes through zero.
With the DS termination the $1\um$ valve gives $0.009$ to $0.022\AW$ at
$300$~K and $0.02$ to $0.09\AW$ at $100$~K, a fifth or less of the funnel
at $300$~K, and a matched antenna does not bring its efficiency at the
drift limit above $4\times10^{-5}$. The shape-only values of this valve
change by a factor of two or more between meshes at $1\THz$ and by up to
a quarter at $0.1$ and $0.3\THz$; they are bounds of the right order. The
valve's asymmetry is an inertial loss that needs a Reynolds number of
ten; at terahertz frequencies the long path that produces it is also
several plasmon wavelengths long, so its stages do not act coherently.

With ideal contacts and the Johnson noise of $\Rdc$ alone, the
noise-equivalent power of the $1\um$ funnel between $0.3$ and $1\THz$ is
$22$ to $28\pWrtHz$ at $100$~K (if the viscosity does not fall with
frequency, $32$ to $60\pWrtHz$ up to $0.9\THz$ and $100\pWrtHz$ at
$1\THz$) and $97$ to $102\pWrtHz$ at $300$~K, or $55$ to $99\pWrtHz$
with leads of $1\um$. That of the $175\nm$ funnel at $100$~K is $6$ to
$29\pWrtHz$ in the same band and $4\pWrtHz$ at its taper resonance.

Table~\ref{tab:sens} lists what the choices of the model do to these
numbers. A contact that fixes the density is the condition of Dyakonov
and Shur and of
Refs.~\cite{tomadin2013theory,svintsov2013hydrodynamic}; one that fixes
the energy carried per particle of the drifting gas fixes
$\Phi(n)+m^{*}u^{2}/4$ instead. The response is linear in the coefficient
of $u^{2}$ in the contact condition, so other choices can be read off
these two. The coefficients $-\tfrac12$ and $-\tfrac14$ of
Eq.~\eqref{eq:mom} hold at any temperature for a flow that exchanges no
heat; for isothermal flow at the reference density they are $-0.34$ and
$-0.17$ at $300$~K and $-0.47$ and $-0.24$ at $100$~K (Supplemental
Material). With the cooling time of $2.5$~ps measured in encapsulated
graphene at room temperature~\cite{tielrooij2018out} the product of
angular frequency and cooling time is $1.6$ at $0.1\THz$ and $16$ at
$1\THz$, so the flow is nearer the first case, which is the one we use.
The viscosity of the $100$~K set is an assumption in its magnitude and
in its dependence on frequency (Table~\ref{tab:params}). In the
viscoelastic form of Sec.~\ref{sec:model-plasmonic} it falls above
$\omega\tauee=1$, at $100$~K above $0.16\THz$. With the viscosity held
at its dc value instead [dotted curves in Fig.~\ref{fig:small}(b)] the
DS response of the $1\um$ funnel at $100$~K is $0.09$ to $0.18\AW$
between $0.1$ and $0.9\THz$ and $0.055\AW$ at $1\THz$, where the
viscoelastic form gives $0.24\AW$, and the share of the shape between
$0.3$ and $1\THz$ is $0.09$ to $0.41$ in place of $0.27$ to $0.38$. The
DS numbers of this section are therefore uncertain by about a third
downward from the contact, by a fifth upward from the temperature at
$300$~K, and by some tens of per cent in either direction from the
magnitude of the viscosity at $100$~K; those at $100$~K above $0.1\THz$
are lower by up to a factor of $2.7$ on the frequencies of the sweep, and
of four at $1\THz$, if the viscosity does not fall with frequency. The
shape-only numbers at $100$~K depend on the viscosity by more than that.
What the fluid description itself leaves out is taken up in
Sec.~\ref{sec:discussion}.

\begin{table*}
\caption{What the choices of the model do to the small-signal response
of the $1\um$ funnel with $12\um$ leads between $0.1$ and $1\THz$.
Factors multiply the values of the text. The runs with a fifth of the
viscosity at $100$~K are on a mesh of $40\nm$ and are less converged
than the others. The uniform channel with an ideal source is more
sensitive than the funnel: the second contact condition leaves between
two thirds and a tenth of its response at $100$~K, and the ordinary
convective term raises it $1.7$ to $2.9$ times at $100$~K and $1.6$ to
$1.9$ times at $300$~K.}
\label{tab:sens}
\begin{ruledtabular}
\begin{tabular}{L{0.26\textwidth}L{0.22\textwidth}L{0.19\textwidth}L{0.25\textwidth}}
Choice varied & DS, $100$~K & DS, $300$~K & Shape only \\
\hline
Contact fixes the energy per particle, not the density & $0.65$ to $0.93$ ($175\nm$ funnel: $0.64$ to $1.02$) & $0.65$ to $0.74$ ($175\nm$ funnel: $0.66$ to $0.74$) & $0.6$ to $1.3$ at $100$~K, $0.6$ to $1.8$ at $300$~K \\[3pt]
Isothermal coefficients in Eq.~\eqref{eq:mom} & at most $4\%$ higher & $16$ to $23\%$ higher & \\[3pt]
A fifth of the viscosity & $1.4$, $0.75$ and $0.8$ at $0.1$, $0.3$ and $1\THz$ & changes by $3$ to $10\%$ & at $100$~K: $4.8$ and $0.6$ at $0.1$ and $0.3\THz$, opposite sign at $1\THz$ \\[3pt]
Viscosity constant in frequency & within $6\%$ below $0.1\THz$; $0.09$ to $0.18\AW$ from $0.1$ to $0.9\THz$; $0.055\AW$ at $1\THz$ & efficiency per watt changes by $1\%$ at $0.1$ and $0.3\THz$, $10\%$ at $1\THz$ & at $100$~K: $0.013$ to $0.073\AW$ from $0.2$ to $1\THz$ \\[3pt]
Ordinary convective term in place of Eq.~\eqref{eq:mom} & up to $1.7$ & about $1.7$ from $0.3$ to $1\THz$ & \\[3pt]
Mesh ($20\nm$ against the extrapolated limit) & up to $6\%$ low & $1\%$ & up to $11\%$ low \\
\end{tabular}
\end{ruledtabular}
\end{table*}

\begin{figure*}
\includegraphics[width=\textwidth]{\figpath/fig4_plasmonic_small_signal.pdf}
\caption{Small-signal plasmonic response. (a) Validation on a uniform
$1\um\times10\um$ channel with an ideal source and open drain: dc drain
voltage per squared ac amplitude in units of the inverse of the voltage
unit $V_{u}$ of each gas, which for the classical gas is the gate
overdrive. The lines are Eq.~\eqref{eq:channel} for the classical gas, the
closed form of Dyakonov and Shur~\cite{dyakonov1996detection}, and for the
Dirac gas of Eq.~\eqref{eq:mom}, at half the classical response; the
solver is shown for both without viscosity at $100$~K and for the Dirac
gas with the parameters of the $300$~K set, for which no line is drawn.
(b)--(d) Current
responsivity into a $75\ \Omega$ antenna for the $1\um$ funnel, the $175\nm$
funnel and the classical Tesla valve, with the drain tied to the gate
(circuit asymmetry) and with the symmetric drive that isolates the shape.
Open symbols mark a negative sign. The green band is $0.3$ to $1\THz$.
Dotted, in (b): the two $100$~K curves with the viscosity held at its dc
value, black for the circuit asymmetry and grey for the shape. In
the grey band the shape-only response changes with the mesh by as much as
its value, so that neither magnitude nor sign is established there. In
(b) the $100$~K curves are drawn to $0.91\THz$, the last frequency of the
sweep below $1\THz$; above $1\THz$ they are not converged with the mesh.}
\label{fig:small}
\end{figure*}

\subsection{Large signal: a square law up to the drift limit}
\label{sec:large}

A detector is judged by responsivity and noise, a harvester by conversion
efficiency $\eta=P_{\mathrm{dc}}/\Pav$ at the power the antenna supplies.
For field-effect detectors that efficiency is quadratic in the amplitude
at small signal, turns linear at an amplitude equal to the gate
swing~\cite{gutin2012plasmonic,rudin2014theory,but2014nonlinear}, and has
an absolute maximum of the order of ten per cent, reached below threshold
where the voltage scale is the thermal
one~\cite{kachorovskii2013performance}. A gapless monolayer has no
subthreshold regime, so the voltage scale of the gated fluid is the gate
overdrive. For the Dirac gas a second scale comes first. A plasma wave of
relative density amplitude $\delta n/n_{0}$ carries a drift
$s_{0}\,\delta n/n_{0}$, and $s_{0}=3.6\,\vF$ under our gate, so a
density swing of $28\%$ would already move the carriers at $\vF$, which
Eq.~\eqref{eq:exact} does not allow. In second-order perturbation the
device is a dc Thevenin source with open-circuit voltage
$V_{\mathrm{oc}}=r\,V_{a}^{2}/V_{u}$ and resistance $\Rdc$, where
$V_{u}=m^{*}s_{0}^{2}/e$ is the natural voltage unit of the gated fluid
($1.49$~V here, the gate overdrive up to the quantum-capacitance
correction) and $r$ is a pure number. Into a matched dc load
\begin{equation}
\eta=2r^{2}\,\frac{\Rant}{\Rdc}\,\frac{V_{a}^{2}}{V_{u}^{2}}
=\frac{16\,r^{2}\Rant^{2}}{\Rdc V_{u}^{2}}\,\Pav .
\label{eq:eta}
\end{equation}
For the $1\um$ funnel between $10\GHz$ and $1\THz$ the small-signal
solver gives $r=0.06$ to $0.11$ and $\eta/\Pav=0.9$ to
$2.9$~W$^{-1}$ at $100$~K and $0.40$ to $0.64$~W$^{-1}$ at $300$~K. What
the nonlinear solver decides is how far the law holds and what ends it.

Figure~\ref{fig:large} shows the time-domain results for the DS-terminated
$1\um$ funnel. The time-domain solver has a constant viscosity and was
run on a mesh of $80\nm$ at the walls; with the viscosity held constant
in both, the two solvers agree to the error of that mesh (Supplemental
Material). Without the viscoelastic term the $300$~K efficiencies per
watt are nearly those of Sec.~\ref{sec:small}, and those at $100$~K are
$1.1$, $1.0$ and $0.11$~W$^{-1}$ at $0.1$, $0.3$ and $1\THz$, where the
small-signal solver gives $1.0$, $2.4$ and $2.4$~W$^{-1}$. Each curve is
therefore compared with the square law of its own solver on its own
mesh, and the $100$~K curve at $1\THz$ is indicative only. The
efficiency meets that law to $1\%$ at the smallest amplitude at $0.1$ and
$0.3\THz$ and stays within $7\%$ of it as long as the largest drift speed
in the device is below the drift limit of $0.3\,\vF$
[filled symbols in Fig.~\ref{fig:large}]. That speed is reached at an
available power of $26$, $38$ and $105\ \mu$W at $0.1$, $0.3$ and $1\THz$
at $100$~K and of $100$ to $130\ \mu$W at $300$~K, and the efficiency
there is $1\times10^{-5}$ to $4\times10^{-5}$ at $100$~K and
$4\times10^{-5}$ to $6\times10^{-5}$ at $300$~K; on a mesh of $40\nm$ at
$0.3\THz$ the power at the limit differs by $1\%$ and the efficiency by
$4\%$ or less. The small-signal solver, with the viscoelastic term, puts
the same speed at $17$ to $22\ \mu$W at $100$~K, with an efficiency at
the limit of $2\times10^{-5}$, $4\times10^{-5}$ and $5\times10^{-5}$ at
the three frequencies in place of $3\times10^{-5}$, $4\times10^{-5}$ and
$1.3\times10^{-5}$. At $300$~K the two solvers agree on the drift to
$3\%$, and the small-signal solver gives $5\times10^{-5}$ to
$7\times10^{-5}$ at the limit. Where the drift is largest depends on the
temperature (Supplemental Material). At $100$~K it is in the neck. At
$300$~K the plasmon is damped along the $12\um$ lead and the drift is
largest at the source contact, across the width of the lead: the powers
at the limit at $300$~K are those of a $4\um$ lead on a $75\ \Omega$
source and not a property of the funnel.

Beyond the limit the response of the ideal gas rises above the square
law, to between $1.1$ and $1.5$ times at $0.8\,\vF$ and between $1.1$ and
$2.5$ times above $0.9\,\vF$ [Fig.~\ref{fig:large}(b)]. A control run
with the second-order terms of Eq.~\eqref{eq:mom} alone stays on the law,
at $0.95$ of it where the full equation gives $1.44$, so the rise comes
from the terms of Eq.~\eqref{eq:exact} beyond second order. At the next
amplitude the drift would exceed $\vF$ and there is no periodic state:
the runs end at $0.33$ to $0.59$~mW at $100$~K and $0.92$~mW at $300$~K,
an amplitude of $0.3$ to $0.5\,V_{u}$, well short of the gate overdrive
(the $100$~K ladder at $1\THz$ was stopped earlier, at $0.56\,\vF$). The
largest efficiency reached is $0.09\%$, at $0.94$ to $0.96\,\vF$, on a
curve that is still rising steeply, so the value as the drift approaches
$\vF$ is not bracketed. The range above $0.3\,\vF$ is that of the ideal
gas and not a prediction for a device: heating and the emission of
optical phonons, which the model lacks, limit the drift of a real sample
inside it. Returning the dc current through the antenna instead of a
matched load costs the expected $15\%$ up to the drift limit and up to
$24\%$ above it.

The shape-only response does not gain on that of the terminations at
large signal. With the symmetric drive the funnel gives $0.19$~W$^{-1}$
at $0.3\THz$ and $100$~K, a fifth of the efficiency of the DS term of the
same device in the same solver; it is $7\%$ below its square law at the
drift limit and falls to $0.77$ of it at $0.54\,\vF$ and $0.45$ at
$0.79\,\vF$ [Fig.~\ref{fig:large}(b)], while the response of the
terminations rises. The picture of a geometric diode that improves when
driven hard is not what the fluid does at $0.3\THz$.

The scale of the result is set by Eq.~\eqref{eq:eta}. A beam of
$1$~W/m$^{2}$ delivers about $10$~nW to a half-wave antenna at $1\THz$,
which the element converts at $\eta\sim10^{-8}$; the tens of microwatts
that the drift limit allows on one element are converted at a few times
$10^{-5}$.

\begin{figure*}
\includegraphics[width=\textwidth]{\figpath/fig5_large_signal.pdf}
\caption{Finite-amplitude response of the DS-terminated $1\um$ funnel from
the time-domain solver with Eq.~\eqref{eq:exact}. Filled symbols: largest
drift speed in the device below $0.3\,\vF$; open symbols: above it. Each
curve ends at the last amplitude for which a periodic state exists, except
that at $1\THz$ and $100$~K, which was stopped at $0.56\,\vF$. (a)
Conversion efficiency into a matched dc load against available power; the
dotted lines are square laws with the coefficients given. (b)
Efficiency relative to the square law of the same solver,
Eq.~\eqref{eq:eta} with its small-signal $r$, normalised at the smallest
amplitude, against the
largest drift speed; the green curve is the shape-only response (symmetric
drive) at $0.3\THz$, $100$~K. (c) Largest drift speed against the ac
amplitude.}
\label{fig:large}
\end{figure*}

\subsection{Operating point and contacts: where the limit lies and what is converted there}
\label{sec:operating}

Equation~\eqref{eq:eta} suggests a lever. The drift limit is reached at an
amplitude of order
$V_{u}\vF/s_{0}$, and both $V_{u}\propto n_{0}d$ and $s_{0}$ fall with the
density and the gate distance. A lower density or a thinner gate should
therefore bring the limit down to the microwatts an
antenna collects. What that does to the efficiency there depends on the
pure number $2r^{2}\Rant/\Rdc$, which the solver has to supply. We swept
$n_{0}=5\times10^{10}$ to $10^{12}\percm$ and $d=30$, $10$, $5\nm$ with the
scattering times of Table~\ref{tab:params}, flagging points outside the
Fermi-liquid limit $\EF\geq3\kB T$ or the gating limit
$e n_{0}/C_{g}\geq4\kB T/e$, and took from the small-signal solution at
each point the available power $P_{\mathrm{lim}}$ at which the largest
drift speed reaches $0.3\,\vF$ and the efficiency of the square law at
that power (Appendix~\ref{app:knee}).

The density moves the limit as expected [Fig.~\ref{fig:operating}(a)]: from
$17\ \mu$W at the reference point to $12\ \mu$W at
$(5\times10^{11}\percm,10\nm)$ and $0.4\ \mu$W at $(5\times10^{10}\percm,30\nm)$,
with $\eta/\Pav$ rising from $2.4$ to $24$ and $790$~W$^{-1}$ at $100$~K
and $0.3\THz$. A thinner gate at the reference density does not lower it:
the limit is at $17$, $37$ and $64\ \mu$W for $30$, $10$ and $5\nm$. The
efficiency at the limit moves less than the limit
[Fig.~\ref{fig:operating}(b)]. Over the
valid grid at $0.3\THz$ it is $4.0\times10^{-5}$ to
$5.4\times10^{-4}$ at $100$~K and $0.6\times10^{-4}$ to $1.1\times10^{-4}$ at
$300$~K, and with the same points at $0.1$ and $1\THz$ included it is
$2\times10^{-5}$ to $5\times10^{-4}$ and $0.4\times10^{-4}$ to
$1.7\times10^{-4}$. A thinner gate raises it, from $4.0\times10^{-5}$ at $30\nm$ to
$4.4\times10^{-4}$ at $5\nm$ at the reference density and $100$~K, although
$r$ falls from $0.10$ to $0.034$: a larger amplitude in units of $V_{u}$,
$0.8$ of the gate overdrive in place of $0.07$, is reached before the drift
limit. A lower density raises $r$, to $0.22$ at the edge of the
Fermi-liquid limit, and the efficiency to $3\times10^{-4}$. A matched
antenna in place of $75\ \Omega$ gives at
best $5.6\times10^{-4}$ at $100$~K and $1.8\times10^{-4}$ at $300$~K. If the mobility rather
than the scattering time is held constant as the density falls, the
low-density gain in $\eta/\Pav$ shrinks by a factor of up to $8.5$ and the
efficiency at the limit is $0.8\times10^{-4}$ to $2.7\times10^{-4}$ below
the reference density. The design rule is
therefore to put the drift limit at the available power; by the square
law the element then converts
at $10^{-5}$ to $5\times10^{-4}$, and nothing in $n_{0}$ or $d$ inside the
valid window raises that number to a tenth of a per cent. These
efficiencies scale as the square of the assumed limit: with $0.2$ or
$0.5\,\vF$ in place of $0.3\,\vF$ they
would be $0.44$ or $2.8$ times as large, the latter if the square law
still held there.

The upper ends of these ranges lie under the thin gates, at amplitudes
of $0.5$ to $1.0$ of the gate overdrive, where the square law of a
field-effect detector is expected to give
way~\cite{gutin2012plasmonic,rudin2014theory,but2014nonlinear}, and they
are values of the small-signal solver, with the viscoelastic term. We
therefore ran the time-domain solver at the five grid points that carry
them (Table~\ref{tab:ladders}). The law holds. All five ladders pass the
drift limit, with the density nowhere below three quarters of its gated
value, and at the limit the efficiency is $1.04$ to $1.12$ times the
square law of the same solver. Two other things lower the numbers. The
first is the drift, which grows faster than the amplitude, so that the
limit is reached at $0.63$ to $0.75$ of the power that the small-signal
solver assigns to it. At $300$~K, where the two solvers are one model,
the efficiency at the limit is for that reason $0.70$ and $0.77$ of the
grid value, $1.2\times10^{-4}$ in place of $1.7\times10^{-4}$ at the top
of the $300$~K range. The second is the viscosity, at $100$~K. The
time-domain solver gives $5\times10^{-5}$ to $7\times10^{-5}$ at the
three points where the grid has $4.4\times10^{-4}$ to
$5.4\times10^{-4}$, and most of that factor of seven to eleven is not a
matter of amplitude: with the viscosity held constant the small-signal
solver itself gives $6\times10^{-5}$ to $1.0\times10^{-4}$ at its limit.
Under a thin gate the fall of the viscosity with frequency thus carries
$81$ to $87\%$ of the efficiency per watt at $0.3\THz$, against a little
over half for the reference device. The largest $100$~K efficiencies of
Fig.~\ref{fig:operating}(b), those under the thin gates, therefore stand
on the viscoelastic form of the viscosity, which is an assumption
(Table~\ref{tab:params}) and which no large-signal calculation here
contains; the low-density points under the $30\nm$ gate, which reach
$3\times10^{-4}$, were not examined in this way. At the seven points
followed in time, with a viscosity that does not fall with frequency,
the element converts at $10^{-5}$ to $1.2\times10^{-4}$ at the drift
limit.

\begin{table*}
\caption{Time-domain check of the efficiency at the drift limit under
thin gates ($1\um$ funnel with $12\um$ leads, DS termination,
$75\ \Omega$). Grid: square law of the small-signal solver, which has
the viscoelastic term, at the amplitude where its first-order drift is
$0.3\,\vF$ [Fig.~\ref{fig:operating}(b)]. Constant $\nu$: the same with
the viscosity held at its dc value. Time domain: efficiency where the
largest drift reaches $0.3\,\vF$, the amplitude there over the gate
overdrive $U_{0}$, and the efficiency over the square law of that
solver, $S$. Efficiencies in units of $10^{-4}$.}
\label{tab:ladders}
\begin{ruledtabular}
\begin{tabular}{cccccccc}
Set & $n_{0}$ (cm$^{-2}$), $d$ (nm) & $f$ (THz) & Grid & Constant $\nu$ & Time domain & $V_{a}/U_{0}$ & $S$ \\
\hline
$100$~K & $5\times10^{11}$, 5 & 0.3 & 5.1 & 0.63 & 0.47 & 0.75 & 1.12 \\
 & $2\times10^{11}$, 10 & 0.3 & 5.4 & 0.96 & 0.70 & 0.47 & 1.04 \\
 & $10^{12}$, 5 & 0.3 & 4.4 & 0.87 & 0.63 & 0.71 & 1.04 \\
$300$~K & $5\times10^{11}$, 5 & 0.1 & 1.7 & & 1.19 & 0.77 & 1.11 \\
 & $10^{12}$, 5 & 0.3 & 0.94 & & 0.72 & 0.71 & 1.12 \\
\end{tabular}
\end{ruledtabular}
\end{table*}

Contact resistance is the dominant practical loss, and it is not a series
loss. A contact resistance $R_{c}$ at each contact of the reference device
reduces $\eta/\Pav$ by factors $2.1$, $7.8$ and $180$ for $R_{c}=100$, $300$
and $1000\ \Omega\um$ at $100$~K and $0.3\THz$
[Fig.~\ref{fig:operating}(c)], the responsivity from $0.24\AW$ to $0.14$,
$0.06$ and $0.01\AW$, and it raises the NEP from $24$ to $35$, $66$ and
$320\pWrtHz$. The series picture underestimates this. The gate carries
displacement current, so the ac drain current is not the negative of the
source current, and a resistance at the drain lifts it off the gate
potential at the signal frequency: the boundary condition that makes the
Dyakonov-Shur mechanism work is partly removed. Matching the antenna to the
new impedance recovers little. That a lower access resistance raises the
photoresponse of graphene detectors strongly has been
measured~\cite{castello2024impact}, and that the shunt
between gate and drain at the signal frequency is what makes a transistor
mix is the basis of the distributed self-mixing
picture~\cite{lisauskas2009rational}; the present result connects the two.

For detection the same lever is worth having
[Fig.~\ref{fig:operating}(d)]. The Johnson-limited NEP of the $1\um$ funnel
at $0.3\THz$ falls from $24\pWrtHz$ at the reference point to
$1.3\pWrtHz$ at $5\times10^{10}\percm$ at $100$~K if the $5$~ps scattering
time survives there, $3.8\pWrtHz$ with constant mobility and $2.7\pWrtHz$ with
$300\ \Omega\um$ contacts; at $300$~K the best valid point is
$27\pWrtHz$. The low-density points are also the least degenerate, and
there the choice between adiabatic and isothermal flow matters most: with
the isothermal coefficients the response of the three least degenerate
valid $100$~K points at $30\nm$ is $6$ to $35\%$ higher and that at
$5\times10^{11}\percm$ and $300$~K is $24$ to $52\%$ higher
(Supplemental Material). The power at the drift limit is the same in
either case.

\begin{figure*}
\includegraphics[width=\textwidth]{\figpath/fig6_operating_point.pdf}
\caption{Carrier density, gate distance and contacts for the DS-terminated
$1\um$ funnel. (a) Available power at which the largest drift speed in the
device reaches $0.3\,\vF$, at $0.3\THz$. (b) Efficiency of the square law
at that power. (c) Efficiency per watt against contact
resistance, relative to ideal contacts, at the reference point. (d)
Johnson-limited NEP at $0.3\THz$. Open symbols are outside the
Fermi-liquid or gating limits.}
\label{fig:operating}
\end{figure*}

\section{The ballistic element at terahertz frequencies}
\label{sec:ballistic}

One mechanism is outside the fluid model by construction. In the clean
limit the ballistic funnel of Sec.~\ref{sec:dc} has a curvature
$G_{2}=3.8$~mA/V$^{2}$; if it responded instantaneously, a cosine drive of
$0.1$~V across the device would rectify $18\ \mu$A ($19\ \mu$A from the
curvature, less the quartic term of the static curve), which is $0.87\AW$ of
the available power of a $75\ \Omega$ source, three to sixteen times the plasmonic
figures. Whether it does respond is a transit-time question. It has been
asked before of a diode of about $100\nm$, with a time-dependent particle
simulation that found rectification into the terahertz range and none
above about $5\THz$~\cite{stearns2020high}; we ask it of the device with
the documented dc asymmetry, then of its length, and for a $0.55\um$
device we compute the field as well.

The imposed potential is made time dependent, $U(x,t)=-eV_{a}\cos(\omega
t)\,g(x)$, and the flights integrate $\hbar\dot{\bm{k}}=-\nabla U$ without
the energy conservation that holds only in a static landscape. Two profiles
$g$ bracket the electrostatics, which the model does not compute: the
uniform ramp over the $6\um$ device used at dc, and a drop of length
$\ell_{F}$ centred on the neck, with $\ell_{F}=175$, $500$ and $1600\nm$
(one neck width, three, and the taper length). The bias-window estimator of
Eq.~\eqref{eq:landauer} relies on flux conservation in a fixed landscape and
fails in a moving one, so electrons are injected from both contacts with the
equilibrium distribution $N(E)f(E-\EF)$ at uniformly distributed phases, and
\begin{equation}
\Idc=\frac{g e}{h}\,\Phi\,\bigl[\langle T_{SD}\rangle-\langle T_{DS}\rangle\bigr],
\label{eq:twosided}
\end{equation}
with $\Phi=\int dE\,N(E)f(E-\EF)$ and the averages taken over injection
phase; $(ge/h)\Phi$ is $1.59$~mA
for the $2\um$ contacts at $300$~K. The first harmonic of the drain current, from the
arrival phases, gives the transport admittance and hence the mismatch to
the antenna. Six further tests validate the estimator: it reproduces the
bias-window current in a static field, it gives zero at zero bias, zero for
the straight control at every frequency, and at $300$~MHz it reproduces the
rectified current and the admittance computed from the static curve
(Supplemental Material). The results are reported as the transfer
$S_{B}(f)=\Idc(f)/\Idc(0)$, with $\Idc(0)$ the quasi-static value for the
same field profile, and as responsivity in two forms. $\RI$ is the current
of one device per available power of the $75\ \Omega$ source, with the
mismatch taken from the first-harmonic admittance. The ballistic devices
have resistances of $0.6$ to $4$~k$\Omega$, far from the antenna, so we
also give the current per watt absorbed at the static conductance $G$,
\begin{equation}
R_{\mathrm{abs}}=\frac{2\,\Idc}{V_{a}^{2}G},
\label{eq:rabs}
\end{equation}
which is what $N=1/(\Rant G)$ devices in parallel would give on a matched
antenna if their impedance did not change with frequency. The plasmonic
funnel of Sec.~\ref{sec:plasmonic} with the DS termination absorbs $87$ to
$100\%$ of the available
power above $0.1\THz$, so its $\RI$ is to be compared with
$R_{\mathrm{abs}}$. We do not normalise by the power that the model itself
absorbs at the drive frequency: with an imposed field and no displacement
current its conduction admittance falls, vanishes and changes sign with
frequency, and a responsivity so defined diverges where it does. For
the one device whose field is computed (Sec.~\ref{sec:pic}) the absorbed
power is a result, and $R_{\mathrm{abs}}$ there is the rectified current
over it.

\subsection{The $6\um$ funnel in a uniform field: two roll-offs}

With the uniform field the response dies in two steps
[Fig.~\ref{fig:ballistic-thz}(a)]. The transfer is $0.98\pm0.04$ at
$0.3\GHz$ and already $0.82$ at $3\GHz$ and $0.67$ at $10\GHz$, although a
direct flight from the source contact to the neck takes $3.6$~ps. The
reason is in the distribution of flight times
[Fig.~\ref{fig:ballistic-thz}(c)]. The taper scatters part of the incoming
electrons into directions nearly perpendicular to the axis, and these drift
back to the source contact, bouncing between the walls of the lead, over
tens of picoseconds: $5\%$ of the source-injected electrons fly longer than
$20$~ps and $0.9\%$ longer than $50$~ps. In a field that fills the leads,
this slow population is displaced in energy between wall collisions once
the period approaches its lifetime, and the transmission in both directions
changes. The response time of a ballistic device is set by its slowest
population, not by its mean transit, and a converging taper manufactures
slow electrons. The second step is the transit roll-off of the direct
electrons: the transfer is $0.41$ at $100\GHz$, crosses zero at $300\GHz$,
where $\sin(\pi f\tau_{\mathrm{tr}})$ vanishes for
$\tau_{\mathrm{tr}}=3.6$~ps, and is $0.05\pm0.04$ at $1\THz$, i.e.
$0.05\AW$. At a mean free path of $700\nm$ the quasi-static responsivity is
only $0.15\AW$ to begin with; it is $0.04\AW$ at $30\GHz$ and zero at
$1\THz$. In a uniform field, the assumption that gives this device its
large computed dc asymmetry, the $6\um$ funnel is not
a terahertz rectifier.

\subsection{Field confined to the neck}

Self-consistent electrostatics concentrates the field of such a diode at
the neck rather than spreading it along the
leads~\cite{truccolo2024comprehensive}. Confining the drop to the
neck lowers the static asymmetry, from $1.61$ for the uniform field to
$1.52$, $1.50$ and $1.40$ for $\ell_{F}=1600$, $500$ and $175\nm$, and it
changes the dynamics [Fig.~\ref{fig:ballistic-thz}(a) and
\ref{fig:ballistic-thz}(b)]. For $\ell_{F}=500\nm$ the transfer is
$1.02\pm0.04$ up to $3\GHz$, since the slow electrons in the lead no
longer see a field, then falls to $0.54$ at $100\GHz$, and is $0.64$ at
$300\GHz$ and $0.38\pm0.04$ at $1\THz$. The decline between $10$ and
$100\GHz$ is not what the $0.5$~ps crossing of the field region would
suggest; the transfer is set by the whole flight between the contacts. At
$1\THz$ the transfer is $0.51$ for a drop of one neck width, $0.38$ for
$500\nm$ and zero for $1600\nm$. The corresponding responsivities of one device are
$0.45\pm0.04$ and $0.37\pm0.04\AW$, one and a half to twice the
Dyakonov-Shur response of the plasmonic funnel at $100$~K; per absorbed
watt the $500\nm$ drop gives $0.72\pm0.07\AW$, three times that response. The
amplitude dependence is that of the dc curve: at $300\GHz$ the rectified
current follows $V_{a}^{2}$ to $0.1$~V and is $0.40$ and $0.24$ of the
square law at $0.2$ and $0.3$~V.

Two further cases show that the terahertz response of the confined field is
not simply the static asymmetry carried up in frequency. At $\lambda=700\nm$
the static asymmetry of the confined-field models is reversed ($A=0.94$ to
$0.97$), the rectified current passes through zero between $30$ and
$100\GHz$, and it returns with the sign of the clean funnel, $0.20\AW$ at
$300\GHz$ and $0.24\AW$ at $1\THz$ ($0.66$ and $0.75\AW$ per absorbed
watt). With half-diffuse walls ($p=0.5$) in the
clean limit the static asymmetry of the $500\nm$ drop is $1.01$, and the
responsivity at $1\THz$ is nevertheless $0.45\pm0.04\AW$ ($0.31\AW$ for the
neck-width drop). In this imposed-field model there is, then, a transit-time rectification
of the funnel with a neck-confined field that survives half-diffuse walls
and bulk scattering, of $0.2$ to $0.45\AW$ for one device; for the walls
of a short funnel the computed field of Sec.~\ref{sec:pic} will give the
opposite answer. It does not
depend on the treatment of the few electrons that reach the Dirac point:
replacing their removal by a classical turning point changes no result by
more than one standard error. It does depend on the confinement of the
field, which this model assumes: at $1\THz$ and
$\lambda=700\nm$ the $6\um$ device gives $0.00\pm0.04\AW$ per absorbed
watt with the uniform field and $0.75\pm0.05\AW$ with the $500\nm$ drop.
The static field computed for this device in Sec.~\ref{sec:dc} is nearer
the second case than the first: $55\%$ of the bias falls on the $500\nm$
at the neck and a quarter along the two leads. Its response at $1\THz$
with a computed field was not calculated, and the field under an ac drive
need not have the static profile. Which of the two figures applies to
the $6\um$ device is therefore open, and the statements below about its
length are statements about a field spread over the device.

\begin{figure*}
\includegraphics[width=\textwidth]{\figpath/fig7_ballistic_thz.pdf}
\caption{The $6\um$ ballistic funnel under an ac drive of amplitude $0.1$~V. (a)
Rectified current relative to its quasi-static value for the uniform field
(clean limit and $\lambda=700\nm$) and for a potential drop of $500\nm$ at
the neck; the dotted line is the transit factor for a flight of $3.6$~ps
from the source contact to the neck. (b) The same ratio at $0.3$ and
$1\THz$ against the length of the drop; open symbols: uniform field. (c)
Fraction of electrons with a flight time longer than $t$ in the static
device at $0.1$~V. Error bars are one standard error.}
\label{fig:ballistic-thz}
\end{figure*}

\subsection{Length scaling}
\label{sec:length}

The two roll-offs of the $6\um$ device are properties of its size. The
model has no length of its own (Sec.~\ref{sec:model-ballistic}): a copy of
the device at half scale has half the current at every bias, and under a
drive at twice the frequency it has half the rectified current. Both
statements were tested, as was the equivalence with scattering at fixed
$\lambda/L$: the static current of device D at $0.1$~V and twice that of its
half-scale copy are $188.0$ and $188.6\ \mu$A, and the rectified current
at $10\GHz$ and twice that of the copy at $20\GHz$ are $11.9\pm0.4$ and
$12.1\pm0.5\ \mu$A. The
frequency response of a given shape is therefore a function of the reduced
frequency $fL/\vF$, and the question whether a ballistic funnel rectifies
at $1\THz$ becomes the question how long it is.

Table~\ref{tab:short} and Fig.~\ref{fig:length}(a) define five devices: the
$6\um$ funnel (G0), its half-scale copy (G1), the same taper and neck with
the leads cut to $250\nm$ (G2), a $0.55\um$ funnel with a $100\nm$ neck
(G3), and a $0.50\um$ funnel with the shoulder and neck of the diode
measured at $28\THz$ in Ref.~\cite{zhu2013graphene}, for which we assume a
45 degree taper (G4). Mirror-symmetric versions of G1 to G4 serve as
controls; at $0.1$, $1$ and $10\THz$ and for both field profiles all 22
control points are zero within $2.1$ standard errors. The confined drop is
one neck width long in G1 to G4; for G0 it is the $500\nm$ drop of the
previous subsection.

With a uniform field and no scattering the five transfer functions fall on
one curve [Fig.~\ref{fig:length}(b)]: above $fL/\vF=0.2$ they agree within
about $0.2$, the response is half its quasi-static value near
$fL/\vF=0.4$, and it is zero from about $1.7$ on; for device D it remains
zero within $0.04$ out to $fL/\vF=30$. The frequency of half response is
thus $0.4\,\vF/L$: $67\GHz$ for the $6\um$ device and $0.7\THz$ for the
$0.55\um$ one, and the response is gone at $280\GHz$ and at $3\THz$. The
same scaling with the size of the scatterer over the Fermi velocity was
found for the ratchet current of Dirac electrons on asymmetric
antidots~\cite{ermann2011relativistic}. Below $fL/\vF=0.2$ the curves
differ, and the difference identifies the gigahertz step of the $6\um$
device. Devices G0 and G1, with leads as long as the taper, share a
plateau at $0.67$ between $fL/\vF=0.03$ and $0.3$; G2 and G3, with short
leads, are still at $0.94$ and $0.90$ at $0.06$, and G4 lies between. The slow population
belongs to the leads, as inferred above from the flight times. Scattering
lowers the cut-off when the device is longer than the mean free path and
not otherwise: for G3 at $\lambda=700\nm$ ($\lambda/L=1.3$) the transfer
is that of the clean device, $0.76$ and $0.32$ at $0.3$ and $1\THz$,
while for G2 ($\lambda/L=0.34$) it is $0.41$ against $0.75$ at
$fL/\vF=0.2$. A cut-off near the Drude damping rate, as reported for the
simulated diode of Ref.~\cite{stearns2020high}, and a cut-off at the
flight time are the two limits of the same behaviour.

\begin{table*}
\caption{Funnels used for the length scaling. Shoulder width, neck width
and lead lengths (source, drain) in nm; $L$ is the length between the
contacts. $A$ is the static asymmetry at $0.1$~V with a uniform field for
no scattering and for mean free paths of $700$ and $100\nm$; $R$ is the
resistance at $\lambda=700\nm$ and $N=1/(\Rant G)$ the number of devices in
parallel that matches $75\ \Omega$. G4 has the shoulder and neck of the
diode of Ref.~\cite{zhu2013graphene} with an assumed 45 degree taper; the
others have the 30 degree taper of device D. The runs for G0 are
independent of those of Sec.~\ref{sec:dc}, which gave $1.62\pm0.02$ and
$1.09\pm0.01$.}
\label{tab:short}
\begin{ruledtabular}
\begin{tabular}{llcccccccc}
 & Description & Shoulder/neck & Leads & $L$ ($\mu$m) & $A$, clean & $A$, 700 nm & $A$, 100 nm & $R$ (k$\Omega$) & $N$ \\
\hline
G0 & device D of Ref.~\cite{nguyen2021geometric} & 2000/175 & 2000, 2420 & 6.00 & 1.61 & 1.10 & & 0.95 & 13 \\
G1 & G0 at half scale & 1000/88 & 1000, 1210 & 3.00 & 1.61 & 1.13 & 1.06 & 1.50 & 20 \\
G2 & G0 with short leads & 2000/175 & 250, 250 & 2.08 & 1.91 & 1.55 & 1.33 & 0.77 & 10 \\
G3 & short funnel & 500/100 & 100, 100 & 0.55 & 1.83 & 1.56 & 1.28 & 1.16 & 15 \\
G4 & proportions of Ref.~\cite{zhu2013graphene} & 400/75 & 169, 169 & 0.50 & 1.60 & 1.30 & 1.15 & 1.54 & 21 \\
\end{tabular}
\end{ruledtabular}
\end{table*}

\subsection{A short funnel at $1\THz$ with an imposed field}
\label{sec:short}

For the $0.55\um$ funnel the reduced frequency at $1\THz$ is $0.55$, and
the device is shorter than the room-temperature mean free path. At
$\lambda=700\nm$ and $1\THz$ the imposed-field model gives
\begin{equation}
R_{\mathrm{abs}}=
\begin{cases}
0.70\pm0.02\AW & \text{uniform field},\\
0.71\pm0.03\AW & \text{drop at the neck},
\end{cases}
\label{eq:imposed}
\end{equation}
and $\RI=0.17$ and $0.14\AW$ for one device of $1.2$~k$\Omega$ on the
$75\ \Omega$ antenna [Fig.~\ref{fig:length}(c)]; G4 gives $0.57\pm0.03$ and
$0.70\pm0.02\AW$ per absorbed watt. Two things have changed with respect to
the $6\um$ device. The two profiles, which differed by the whole effect at
$6\um$, agree within the statistical error [Fig.~\ref{fig:length}(d)]. And
the result holds at the mean free path of encapsulated graphene at room
temperature, not only in the clean limit, in which the same device gives
$0.95\pm0.05$ and $0.83\pm0.06\AW$. Taken at face value this is about
three times the response of the plasmonic funnel at $100$~K and thirteen
times its $0.05\AW$ at $300$~K.

In the same model the edges and scattering cost a factor of up to five.
With half-diffuse walls ($p=0.5$) the figures at $1\THz$ and
$\lambda=700\nm$ are $0.42\AW$ (uniform) and $0.14\AW$ (confined); with a
mean free path of $100\nm$ they are $0.22\pm0.04$ and
$0.42\pm0.04\AW$. We also take this case as an indication of what
electron-electron collisions do at $300$~K, where their mean free path
is of that length. It is not a bound: such collisions randomise the
direction of an electron while conserving the momentum of the gas,
which a momentum-relaxing path does not, and they were not simulated. With the confined drop the static asymmetry of G3 falls
below one when scattering is strong ($A=0.81$ at $100\nm$, $0.90$ with
half-diffuse walls at $700\nm$), while at $1\THz$ the current has the sign
of the clean funnel in every case, as for the $6\um$ device. The measured
devices bear the dependence on the mean free path out in order of
magnitude: the diode of Ref.~\cite{zhu2013graphene}, with a measured mean
free path of about $45\nm$, had a zero-bias dc responsivity of $0.12\AW$,
against $0.71\AW$ for G4 at $100\nm$ and $1.33\AW$ at $700\nm$ in the
quasi-static limit of the uniform-field model, and the rectifiers of
chemical vapor deposited graphene of Ref.~\cite{hemmetter2026enhanced},
with mean free paths of at most $65\nm$, have a responsivity of $0.8$~V/W
referred to the beam power, against hundreds of V/W for encapsulated
ballistic rectifiers~\cite{auton2017terahertz}. The amplitude dependence
is a square law with a low ceiling: with the uniform field the rectified
current of G3 follows $V_{a}^{2}$ up to $0.1$~V and is $0.80$ and $0.57$ of
the law at $0.2$ and $0.3$~V, and with the confined drop it is already
$19\%$ below the law at $0.1$~V. The scale is the Fermi energy,
$\EF/e=0.14$~V, a tenth of the gate overdrive and comparable to the
amplitude of $0.1$ to $0.3$~V at which the plasmonic element reaches its
drift limit.

Above $1\THz$ the two field profiles separate again. With the uniform field
the response of G3 and G4 is zero at $3\THz$, as the scaling requires. With
the drop confined to the neck it is $0.85\pm0.03$ and $1.03\pm0.03\AW$ at
$3\THz$ and has fallen by an order of magnitude at $10\THz$, to
$0.05\pm0.03$ and $0.09\pm0.03\AW$ at $\lambda=700\nm$ ($0.23\pm0.04\AW$
for G4 at $100\nm$); at $30\THz$ no device, profile or mean free path gives
a current distinguishable from zero. The response at $28\THz$ measured
for a diode of the size of G4~\cite{zhu2013graphene} is not reproduced. At
that frequency the photon energy, $116$~meV, is close to the Fermi energy
and a classical trajectory in a classical field is not an adequate
description; interband absorption, which the model lacks, is the natural
candidate for what was measured.

The agreement of the two profiles in Eq.~\eqref{eq:imposed} suggests that
for a device this short the field does not matter. The next subsection
tests that by computing it, and finds that it does.

\subsection{The short funnel with the field computed}
\label{sec:pic}

We take the $0.55\um$ funnel G3 into the model of
Sec.~\ref{sec:model-pic}, in which the potential follows from the charge
of the electrons, a gate and metal contacts. The operating point is that
of Eq.~\eqref{eq:imposed}: $\lambda=700\nm$, $1\THz$, an amplitude of
$0.1$~V between the contacts, the drive that leaves only the shape, and
the gate at $300\nm$. The funnel rectifies $0.429\pm0.012\ \mu$A of an
absorbed $3.44\ \mu$W,
\begin{equation}
R_{\mathrm{abs}}=\frac{\Idc}{P_{\mathrm{abs}}}=0.124\pm0.003\AW ,
\label{eq:computed}
\end{equation}
and $\RI=0.023\AW$ for one device on the $75\ \Omega$ antenna, with an
admittance of $0.69-0.17i$~mS. Under the $30\nm$ gate the figures are
$0.132\pm0.003$ and $0.026\AW$, with $0.73+0.34i$~mS. The absorbed power
is now a computed quantity; in the definition of Eq.~\eqref{eq:rabs}, for
comparison, the values are $0.10$ and $0.11\AW$ against $0.70\AW$. The
errors are statistical, from seven and six runs. Changes of the grid, the
time step, the particle number and the length of the contacts move the
current between $0.38$ and $0.49\ \mu$A, and the mirror-symmetric version
of the device, which must give zero under this drive, gives currents of up
to $0.06\ \mu$A in single runs at $0.3$ and $1\THz$ (the mean of six runs
from $0.3$ to $3\THz$ is $0.01\pm0.02\ \mu$A); we take $0.05\ \mu$A as the
floor of the method and quote $0.12\pm0.015\AW$. The computed field leaves a sixth of
what either imposed field gave, or a seventh when both are expressed in
the definition of Eq.~\eqref{eq:rabs}.

The loss is already present at dc. The static asymmetry at $0.1$~V is
$1.053$ under both gates, against $1.56$ and $1.13$ for the two imposed
profiles, while the low-bias conductance, $0.86$~mS ($0.85$~mS from the
currents at $\pm0.1$~V), is that of the uniform-field
model: only the even part of the current-voltage curve has changed.
Figure~\ref{fig:pic}(a) and \ref{fig:pic}(b) show why. About $61\%$ of the
bias drops within $90\nm$ of the neck, centred on it; $14\%$ along the
lead and the taper; and $25\%$ at the two contact planes, $6$ to $9\%$ at
the source and $17$ to $19\%$ at the drain, where it contributes nothing
to the asymmetry. The drop at the source is of the size expected for a
ballistic contact, half of the current over the Sharvin conductance of
the $500\nm$ lead, or $8\%$ of the bias. Neither imposed profile
resembles this. The uniform ramp puts most of the bias along the taper,
which is where its large asymmetry comes from, and the confined drop puts
all of it at the neck: with $61\%$ there and a rectified current that goes
as the square of the bias, the confined-drop model would give
$0.9\ \mu$A in the quasi-static limit, and the computed curve gives
$1.07\pm0.02\ \mu$A.

The frequency dependence is that of the uniform-field model at a seventh
to a ninth of its height [Fig.~\ref{fig:pic}(c)]. At $0.1\THz$ the
rectified current is $1.04\pm0.01\ \mu$A, the quasi-static value of the
computed static curve; it is $0.94$, $0.71$, $0.43$ and $0.15\ \mu$A at
$0.3$, $0.6$, $1$ and $1.8\THz$, half at $0.9\THz$ or $fL/\vF=0.5$. The
transit roll-off of the rectified current thus survives the computation
of the field, for a device whose contacts lie $100\nm$ from the taper, so
that the length between the contacts and the length of the region that
carries the field cannot be told apart. The agreement of the two imposed profiles at $1\THz$ does
not: it was the crossing of a curve that falls with frequency and one
that rises, and no sign of insensitivity to the field. Above $1\THz$ the
gate matters. Under the $300\nm$ gate the absorbed power falls steadily,
to $0.8\ \mu$W at $3\THz$, with $0.19\pm0.02\ \mu$A rectified there and
nothing at $6\THz$. Because the absorbed power falls as well, the response per absorbed watt
does not follow the current: it is $0.08\pm0.01\AW$ at $1.8\THz$ and
$0.24\pm0.03\AW$ at $3\THz$, twice its value at $1\THz$. Under the $30\nm$ gate the device is a plasmonic
cavity between two contacts that hold the potential, with a half-wave
resonance near $3.6\THz$ and a full-wave one near $7\THz$ for the plasma
speed of $3.9\,\vF$; the antisymmetric shape-only drive couples to the
second, the absorbed power rises to $10.7\ \mu$W at $6\THz$, and the
rectified current stays at $0.35$ to $0.50\ \mu$A up to $3\THz$.

With a gate in the model, shape and circuit can be compared in the
ballistic gas as they were in the fluid [Figs.~\ref{fig:pic}(d) and
\ref{fig:pic}(e)]. The circuit term is the current of a straight strip
of the same length and lead width with its drain tied to the gate, the
Dyakonov-Shur termination, here without collisions. Under the $300\nm$
gate it is $0.13$, $0.34$ and $0.01\ \mu$A at $0.3$, $1$ and $3\THz$,
against $0.94$, $0.43$ and $0.19\ \mu$A for the shape term: the shape is
the larger at each frequency. Under the $30\nm$ gate the strip gives
$0.67$, $1.15$ and $-1.77\ \mu$A against $0.89$, $0.48$ and $0.50\ \mu$A.
At $1\THz$ the circuit term is then more than twice the shape term,
$0.20\AW$ per absorbed watt against $0.13\AW$, and the funnel with its
drain tied to the gate, which has both, gives $1.01\ \mu$A, no more than
the strip: the two terms do not add. At $0.3\THz$ the shape term is still
the larger of the two. At $3\THz$, near the half-wave resonance that this
termination does excite, the funnel absorbs four times its
low-frequency power and the strip's current is reversed; the run was
repeated with a longer settling time and did not change, and we have not
looked for the cause. With the gate closer than the device is long,
then, the gate connection rectifies more than the shape at $1\THz$ in
the ballistic gas as well.

How much of the shape term survives in a real device is again a matter of
scattering and of the edges, and here the computed field changes the
answer of the imposed one in both directions [Fig.~\ref{fig:pic}(f)].
The mean free path matters little. Per absorbed watt at $1\THz$ the
funnel gives $0.130\pm0.005\AW$ with no scattering, $0.124$ at $700\nm$
and $0.107\pm0.007\AW$ at $100\nm$ ($0.145$, $0.132$ and $0.099\AW$ under
the $30\nm$ gate), where the uniform-field model lost a factor of three.
Scattering does lower the static asymmetry, $1.09$, $1.05$ and $1.02$,
and the quasi-static current, $2.1$, $1.1$ and $0.23\ \mu$A, but the
current at $1\THz$ falls much less, $0.50$, $0.43$ and $0.22\ \mu$A: what
a long mean free path adds to the static asymmetry is not available at
$1\THz$, and at $100\nm$ the roll-off has moved above $1\THz$. The
conductance halves at $100\nm$, so one device on the antenna does lose a
factor of two, $\RI=0.012\AW$. The edges decide. With half-diffuse walls
($p=0.5$, $\lambda=700\nm$) the static asymmetry is $1.005$ and the
current at $1\THz$ is $-0.07\pm0.01\ \mu$A under the $300\nm$ gate and
$-0.13\pm0.01\ \mu$A under the $30\nm$ gate, a sixth to a quarter of the
specular current with the opposite sign; two independent runs agree, and
the mirror-symmetric control with the same walls gives zero within the
floor. The imposed fields had given $0.42$ and $0.14\AW$ for this case;
the computed field gives at most $0.04\AW$ in magnitude. The forward
shape rectification of the short funnel does not survive walls that
scatter half of all electrons diffusely; what replaces it is a weaker
current of the opposite sign, whose mechanism we have not identified.

How specular the walls must be is shown in Fig.~\ref{fig:pic}(g), which
plots the response per absorbed watt. With
one specularity for all angles, between $p=1$ and
$p=0.5$ the current at $1\THz$ falls almost linearly, by $0.10$ to
$0.12\ \mu$A for each tenth of the reflections that is made diffuse. It
is $0.29\pm0.01$, $0.16\pm0.02$ and $0.10\pm0.02\ \mu$A at $p=0.9$,
$0.8$ and $0.7$ under the $300\nm$ gate and $0.35\pm0.02$,
$0.17\pm0.02$ and $0.08\pm0.01\ \mu$A under the $30\nm$ gate; the
values at $p=0.7$ are no more than twice the floor. The absorbed power
changes by $6\%$ over the scan, so the response per absorbed watt follows
the current: $0.09$, $0.05$ and $0.03\AW$ under the $300\nm$ gate and
$0.10$, $0.05$ and $0.02\AW$ under the $30\nm$ gate. One diffuse
reflection in ten
costs $30\%$ of the response, half is left at $p=0.84$, and the current
passes through zero at $p=0.58$ and $0.63$ for the two gates, by linear
interpolation.

One specularity for all angles is not how a rough edge reflects, and for
this device the difference is large. The walls of the taper are inclined
at $30^{\circ}$ to the axis, so that an electron travelling along the axis
meets them at $60^{\circ}$ from the normal, where Eq.~\eqref{eq:rough}
gives the fourth root of the head-on specularity. With walls that follow
Eq.~\eqref{eq:rough} the current at $1\THz$ under the $300\nm$ gate is
$0.34\pm0.02$, $0.31\pm0.02$, $0.21\pm0.02$, $0.14\pm0.01$,
$0.04\pm0.02$ and $-0.10\pm0.02\ \mu$A for $p_{0}=0.81$, $0.74$, $0.49$,
$0.36$, $0.21$ and $0.05$, that is, for $h=1.1$, $1.3$, $2.0$, $2.4$,
$3.0$ and $4.1\nm$; under the $30\nm$ gate it is $0.37\pm0.02$,
$0.38\pm0.01$, $0.20\pm0.02$, $0.14\pm0.01$, $-0.02\pm0.01$ and
$-0.18\pm0.02\ \mu$A [Fig.~\ref{fig:pic}(g)]. Per absorbed watt the
funnel gives $0.10\AW$ at $p_{0}=0.81$, $0.06\AW$ at $0.49$ and
$0.04\AW$ at $0.36$. The response is half its specular value at
$p_{0}=0.50$ and $0.55$ for the two gates and passes through zero at
$0.16$ and $0.23$, by linear interpolation, where one specularity for
all angles required $0.84$ and about $0.6$. Placed on the scan of the
single specularity, the five rough edges with a positive current under
the $300\nm$ gate rectify like walls of $p=0.94$, $0.91$, $0.84$, $0.77$
and $0.63$; Eq.~\eqref{eq:rough} at $60^{\circ}$ gives $0.95$, $0.93$,
$0.84$, $0.78$ and $0.67$, and averaged over the cosine-law spread of
hits of a gas at rest it gives $0.87$, $0.82$, $0.64$, $0.53$ and
$0.39$. The oblique hits decide. The gas as a whole does not favour
them: counted over all electrons, the fraction of specular reflections
is the same on the sloping walls as on the others to within $0.005$. We
attribute the result to the driven part of the distribution, which
carries the rectified current and moves along the axis; the simulation
does not identify those electrons separately. At $p_{0}=0.05$ the
current is reversed, as it is for half-diffuse walls of a single
specularity. The reversed current is not small: at $-0.10$ and
$-0.18\ \mu$A it is two to nearly four times the floor and $0.24$ and
$0.38$ of the specular current, or $0.03$ and $0.05\AW$ of the opposite
sign. No rougher edge was computed.

Measured specularities of etched graphene edges span this range.
Transverse magnetic focusing, in which the ratio of the second focusing
peak to the first measures the reflection of electrons that meet the
edge near normal incidence, gives $0.2$ to $0.5$ for graphene on
hexagonal boron nitride with edges etched in an oxygen
plasma~\cite{taychatanapat2013electrically}, and $0.8$ and $0.65$ at two
carrier densities for an encapsulated twisted
bilayer~\cite{berdyugin2020minibands}. Collimated focusing in an etched
encapsulated device, fitted by ray tracing, gives fully diffuse
reflection for electrons and two diffuse reflections in three for
holes~\cite{barnard2017absorptive}, and the magnetoresistance of etched
wires points to a specularity near zero~\cite{masubuchi2012boundary}.
Read as head-on values, these place a funnel with etched edges anywhere
between no shape rectification and four fifths of the specular
response: an edge with $p_{0}$ of $0.5$ to $0.8$ keeps $0.06$ to
$0.10\AW$, an edge with $p_{0}$ near $0.2$ keeps nothing, and a rougher
one gives the weaker reversed current. The
comparison has four weaknesses. The measured values are taken as they
were reported, at the carrier densities of the experiments and without
the rescaling to the Fermi wave vector of this device that
Eq.~\eqref{eq:rough} would call for, and the upper end of the range comes
from a twisted bilayer. Equation~\eqref{eq:rough} is the
Kirchhoff result for a surface with Gaussian height statistics; it
ignores the correlation length of the roughness and is least reliable
towards grazing incidence~\cite{shao2017probing}, where it makes the
edge most specular. No measurement on graphene that we know of resolves
the angle of incidence, so the angle dependence is untested there. And
the wire and collimator experiments do not select normal incidence, so
that their low values imply edges at least as rough as a head-on reading
of them suggests.

The contacts are the remaining idealisation, and a quarter of the bias
falls at them even when they are ideal. A contact that reflects might be
expected to act as a resistor in series. For a contact of transparency
$T$ that resistance is $(1-T)/T$ times the Sharvin resistance of the
$500\nm$ lead, which leaves a share $\beta$ of the applied voltage on the
device, $0.90$, $0.75$ and $0.50$ for $T=0.75$, $0.5$ and $0.25$; a
square-law rectifier behind it gives $\beta^{3}$ of its current, since the
rectified current is divided by the series resistance as well, and
$\beta^{2}$ of its response per absorbed watt. The funnel loses much less
[Figs.~\ref{fig:contacts}(a) and \ref{fig:contacts}(b)]. With one
transparency for all angles the current at $1\THz$ under the $300\nm$
gate is $0.35\pm0.02$, $0.36\pm0.01$ and $0.27\pm0.01\ \mu$A at $T=0.75$,
$0.5$ and $0.25$, against $0.43\ \mu$A between ideal contacts and $0.31$,
$0.18$ and $0.06\ \mu$A for the series resistor. The absorbed power
falls in the same proportion, from $3.4$ to $2.2\ \mu$W, so the response
per absorbed watt does not change: $0.107\pm0.005$, $0.125\pm0.005$ and
$0.120\pm0.003\AW$. Under the $30\nm$ gate it is $0.123\pm0.004$,
$0.115\pm0.004$ and $0.071\pm0.005\AW$, where the series resistor would
leave $0.108$, $0.075$ and $0.034\AW$. The values at $T=0.5$ and $0.25$
are means of two or three runs with different seeds, settling times or
particle numbers (Supplemental Material), which agree within $0.4$
standard errors; the mirror-symmetric funnel between contacts of $T=0.25$ gives
$0.01\pm0.01$ and $-0.03\pm0.02\ \mu$A. One device on the antenna loses
more, because its conductance falls: $\RI$ is $0.019$, $0.020$ and
$0.015\AW$ under the $300\nm$ gate.

The static curve shows the same departure from a series connection
[Fig.~\ref{fig:contacts}(c)]. At $T=0.5$ and $0.25$ the two contact
planes take $43$ to $44\%$ and $61$ to $63\%$ of the bias, close to the
$43\%$ and $62\%$ of the series picture, yet the conductance is $0.70$
and $0.44$ of its value between ideal contacts where the series picture
gives $0.75$ and $0.51$, and the asymmetry is $1.065$ and $1.057$,
above the $1.053$ of ideal contacts, where the series picture gives
$1.030$ and $1.014$. The device is shorter than the mean free path, so
an electron that a contact reflects returns to the neck without having
lost its direction. Device and contacts then scatter as one object and
their resistances do not add. We offer this as the likely reason and
have not tested it separately.

What a contact does with electrons that arrive obliquely matters more
than its resistance. With the collimating contact of
Eq.~\eqref{eq:collimate}, at the two junction widths that give a strip
the conductance of $T=0.5$ and $0.25$, the current at $1\THz$ is
$0.08\pm0.02$ and $0.04\pm0.01\ \mu$A under the $300\nm$ gate and
$-0.05\pm0.02$ and $-0.06\pm0.02\ \mu$A under the $30\nm$ gate: at or
near the floor, $11$ to $17$ standard errors below the contacts of the same
conductance that do not select by angle, and no more than $0.03\AW$ in
magnitude. These are single runs in which the gas was still heating
(below), so the size of what is left is uncertain; the loss relative to
contacts that do not select by angle is not. The shape rectification of
the short funnel therefore needs
contacts that absorb at wide angles. Measured contact resistances do not
say whether a contact does. The one-dimensional edge contacts to
encapsulated graphene show about $150\ \Omega\um$ at a density of
$2.5\times10^{12}\percm$~\cite{wang2013one}, a transparency near $0.4$
in the first model and inside the range computed, while scalable edge
contacts at about $2$~k$\Omega\um$~\cite{pandey2018all} correspond to a
transparency below $0.1$ and are outside it. Chromium-gold edge contacts
dope the graphene beside them $n$-type~\cite{handschin2017fabry}, which
puts a junction in front of the contact for a hole channel and none for
an electron channel.

Three reservations apply to the contact runs. The model has no loss of
energy to the lattice, and the driven gas is warmer the less its
contacts absorb: its mean kinetic energy is $2.5\%$ above the
equilibrium value between ideal contacts, $3$, $5$ and $7\%$ at
$T=0.75$, $0.5$ and $0.25$, and $7$ and $14\%$, still rising after
$20$ periods, between the collimating contacts. The two contact models
are the extremes of angular selection that could be stated simply, and
neither is a model of a metal edge. And the distance to the contacts was
not separated from the transit time: with ideal contacts and leads of
$300\nm$ in place of $100\nm$ the current is $0.54$ and $0.67$ of its
value for the two gates, at $fL/\vF=0.95$ in place of $0.55$.

The amplitude dependence differs from that of the imposed fields. At
$0.05$~V the current is $0.10\pm0.02\ \mu$A under both gates, consistent
with the square law through the $0.1$~V point. A quadratic and a quartic
term fitted to the points at $0.1$ and $0.2$~V reproduce it and give a
small-signal response of $0.10\AW$ under the distant gate; the $0.12\AW$ of
Eq.~\eqref{eq:computed} is the value at $0.1$~V, which is $0.7\,\EF/e$. At $0.2$~V, which is above
$\EF/e$, it is $2.56$ and $2.66\ \mu$A, $1.5$ and $1.4$ times the square
law, where both imposed fields fell below it; the absorbed power follows
$V_{a}^{2}$ within $7\%$, so $R_{\mathrm{abs}}$ rises to $0.20\AW$. Into a
matched dc load the conversion efficiency estimated from the
short-circuit current and the static conductance,
$\eta=\Idc^{2}/(4GP_{\mathrm{abs}})$, is $2\times10^{-5}$ at $0.1$~V,
where the device absorbs $3.4\ \mu$W, and $1.5\times10^{-4}$ at $0.2$~V
and $13\ \mu$W. The second figure is the largest amplitude computed and
not a ceiling, since the current there still rises faster than the
square law. As a detector, $N=1/(\Rant\,\mathrm{Re}\,Y_{1})=19$
devices in parallel with their reactance tuned out absorb the available
power, and the Johnson noise of their static conductance gives a
noise-equivalent power of $130\pWrtHz$ at $1\THz$ and $300$~K, against
about $100\pWrtHz$ for the plasmonic funnel at the same temperature.

The outcome for the short funnel is therefore narrower than the imposed
fields suggested. A $0.55\um$ funnel with specular edges rectifies at
$1\THz$ by its shape, at $0.11$ to $0.13\AW$ per absorbed watt under the
distant gate, and $0.10$ to $0.15\AW$ under the near one, for any
mean free path from $100\nm$ up. That is about twice the plasmonic
funnel at $300$~K, $0.05\AW$, and below its $0.20$ to $0.26\AW$ at
$100$~K between $0.3$ and $1\THz$; it is not the factor of thirteen of
Eq.~\eqref{eq:imposed}. The plasmonic figures are per available watt for
one device that absorbs nearly all of it; the ballistic figure is reached
on an antenna only by about nineteen devices in parallel. Rough
edges take from it according to their head-on specularity: half is left
at $0.5$, nothing near $0.2$, and a reversed current of a quarter to two
fifths of the magnitude at $0.05$; etched edges have been measured on
both sides of those marks. Contacts that reflect take little from it as
long as they absorb at every angle, and nearly all of it if they pass only
electrons near normal incidence. Under a close gate the circuit term
exceeds it at $1\THz$.

\begin{figure*}
\includegraphics[width=\textwidth]{\figpath/fig10_ballistic_length.pdf}
\caption{Length scaling of the ballistic response at an ac amplitude of
$0.1$~V. (a) The five funnels of Table~\ref{tab:short} to scale, and the
two short ones enlarged. (b) Rectified current relative to its
quasi-static value against the reduced frequency, for a uniform field and
no scattering; the band marks the fall from half response to zero. (c)
Responsivity of the $0.55\um$ funnel G3 at $\lambda=700\nm$: per absorbed
watt, Eq.~\eqref{eq:rabs} (thick), and for one device on a $75\ \Omega$
antenna (thin, open symbols). (d) Responsivity per absorbed watt at
$1\THz$ against device length for the two field profiles. The blue band in
(c) and (d) is the range of the plasmonic funnel of
Sec.~\ref{sec:plasmonic}, $0.05$ to $0.26\AW$. Error bars are one standard
error.}
\label{fig:length}
\end{figure*}

\begin{figure*}
\includegraphics[width=\textwidth]{\figpath/fig11_selfconsistent.pdf}
\caption{The $0.55\um$ funnel G3 with the field computed; $\lambda=700\nm$
and an amplitude of $0.1$~V except in (f). (a) Static potential energy at
a bias of $0.1$~V with the gate at $300\nm$, in units of the bias;
contours every $0.1$; the dotted rectangles are the contact strips. (b)
The same along the axis for the two gate distances, against the two
imposed profiles of Sec.~\ref{sec:short}; colours and line styles as in
(c). (c) Rectified current against
frequency under the drive that leaves only the shape, for the imposed
fields and for the computed field with the gate at $300\nm$ (full
symbols) and $30\nm$ (open symbols); points within two standard errors of
zero are not drawn. (d), (e) Shape against circuit for the two gate
distances: the funnel under the shape-only drive, a straight strip of the
same length with its drain tied to the gate, the funnel with its drain
tied to the gate, and the mirror-symmetric funnel under the shape-only
drive, which must give zero. (f) Responsivity at $1\THz$ for no
scattering, mean free paths of $700$ and $100\nm$, and half-diffuse
walls: per absorbed watt for the computed field (diamonds, full and open
as in (c)), and by Eq.~\eqref{eq:rabs} for the imposed fields (circles:
uniform; squares: drop at the neck). For the imposed fields the
conductance in Eq.~\eqref{eq:rabs} is here taken from the currents at
$\pm0.1$~V, which places these points up to $5\%$ above the values in
the text. The band is the range of the
plasmonic funnel. (g) Responsivity per absorbed watt at $1\THz$ for the
computed field against the specularity of a head-on hit at the Fermi wave
vector: walls of one specularity $p$ for all angles (grey circles) and
rough walls that follow Eq.~\eqref{eq:rough} (black diamonds); full
symbols and solid lines for the gate at $300\nm$, open symbols and dashed
lines for $30\nm$. The bar marks the range measured for etched
edges~\cite{taychatanapat2013electrically,masubuchi2012boundary,barnard2017absorptive,berdyugin2020minibands}. Error bars are one
standard error.}
\label{fig:pic}
\end{figure*}

\begin{figure*}
\includegraphics[width=\textwidth]{\figpath/fig12_contacts.pdf}
\caption{The $0.55\um$ funnel G3 between contacts that reflect, with the
field computed; $\lambda=700\nm$, specular edges, an amplitude of
$0.1$~V and the drive that leaves only the shape. (a) Rectified current
at $1\THz$ against the transparency $T$ of the two contacts: one
transparency for all angles (black diamonds) and the collimating contact
of Eq.~\eqref{eq:collimate} (orange squares), placed at the transparency
that gives a clean strip the same conductance; full symbols for the gate
at $300\nm$, open symbols for $30\nm$. The upper axis is the contact
resistance that a transfer-length measurement would report. The lines
are a resistor in series with the funnel as it is between ideal contacts
(solid: gate at $300\nm$; dashed: $30\nm$); the grey band is the floor of
the method. (b) The same per absorbed watt. (c) Static potential energy
on the axis at a bias of $0.1$~V, in units of the bias, with the gate at
$300\nm$. Error bars are one standard error.}
\label{fig:contacts}
\end{figure*}

\section{Discussion}
\label{sec:discussion}

\subsection{Shape and circuit, regime by regime}

Table~\ref{tab:verdict} and Fig.~\ref{fig:synthesis} collect the results.
The Tesla valve fails in graphene in every regime, and for one reason. Its
asymmetry is an inertial loss: it needs a fluid, and a Reynolds number of
order ten. The ballistic electron gas is not a fluid, so the loops only add
resistance. The viscous electron liquid is one, but in graphene its
Reynolds number is capped by the Fermi velocity at values where the valve
gives a per cent or less, and the gate adds an asymmetry that is the same
for a straight strip. At terahertz frequencies the valve is long compared
with the plasmon wavelength and its resistance is two orders of magnitude
from an antenna. The one electronic Tesla valve with a large
diodicity~\cite{sarypov2026electron} works in another material, above a
threshold current, in a regime that is not the steady laminar flow of this
paper.

The funnel rectifies in every regime, and what does the rectifying differs.
In the electron fluid it is the circuit. In steady flow the shape leaves an
inertial asymmetry of less than $10^{-3}$ inside the drift limit at the
parameters of Table~\ref{tab:params}, the same with and without the gate,
while the connection of one contact to the gate
produces a per cent at $43$~mV in any channel. At terahertz frequencies the
gated funnel's rectification is that of its terminations: tying the drain
to the gate and driving the source produces $0.1$ to $0.3\AW$ in a funnel,
as it would in a plain channel of the same impedance, while the symmetric
drive that leaves only the shape produces a comparable response only near a
taper resonance at low temperature, and that response saturates first at
large signal.

In the ballistic gas it is the shape, together with the potential that the
bias lays over it, and the first condition for it to work at a frequency
$f$ is one of size:
\begin{equation}
L\lesssim0.4\,\frac{\vF}{f} .
\label{eq:criterion}
\end{equation}
This is the half-response point of Fig.~\ref{fig:length}(b), for a field
spread over the device; with the field computed the $0.55\um$ funnel is at
half response at $0.5\,\vF/L$. At $1\THz$ the criterion gives $0.4\um$.
The $6\um$ funnel on which the dc asymmetry of geometric diodes is best
documented is fifteen times too long, and its response at $1\THz$ exists
only if the potential drop is confined to the neck. A statement that
geometric diodes are or are not terahertz rectifiers has no content
without the length, and the dc asymmetry of a long device is no guide: the
best dc asymmetry at $\lambda=700\nm$ in Table~\ref{tab:short} belongs to
devices between $0.55$ and $2\um$, and only the shortest of them still
responds at $1\THz$ in a uniform field.

Length is necessary and not sufficient. The $0.55\um$ funnel meets the
criterion, and with an imposed potential it appears to give $0.7\AW$ per
absorbed watt whatever the profile. With the potential computed it gives
$0.12\AW$ at the room-temperature mean free path, because the bias does
not go where either imposed profile puts it: a quarter is spent at the
contacts and the rest is centred on the neck. This is a caution beyond the
present device. A uniform field is a common assumption of simulations of
geometric diodes; self-consistent electrostatics was already known to
roughly halve the static asymmetry~\cite{truccolo2024comprehensive}, and
here it reduces a static asymmetry of $1.56$ to $1.05$ and the terahertz
response sixfold. Four further conditions follow from
Sec.~\ref{sec:pic}. The mean free path is the least of them: from no
scattering to $100\nm$ the response per absorbed watt changes from $0.13$
to $0.10\AW$. The edges are the least known. With one specularity for
all angles the response would be halved by one diffuse reflection in six;
with the angle dependence of a rough edge, Eq.~\eqref{eq:rough}, it is
halved at a head-on specularity of $0.5$ and lost at $0.2$, because the
walls of the taper are met obliquely. The values measured for etched
edges, from near zero to
$0.8$~\cite{taychatanapat2013electrically,masubuchi2012boundary,barnard2017absorptive,berdyugin2020minibands},
lie on both sides of those marks. The contacts are the third, and
their resistance is the wrong measure of them. Contacts that reflect
three quarters of the arriving electrons at every angle, at
$340\ \Omega\um$ each, leave $0.12\AW$ per absorbed watt under the far
gate and $0.07\AW$ under the near one, where the same resistance in
series would leave $0.03\AW$; contacts of the same conductance that pass
only electrons near normal incidence leave $0.03\AW$ or less, of
either sign. A
device shorter than the mean free path and its contacts are one
scatterer. And the gate
must be far: under
a gate at $30\nm$ a straight strip with its drain tied to the gate
rectifies twice the current of the funnel shape at $1\THz$, so that the
circuit takes over in the ballistic gas as it did in the fluid.
Shape-based rectification at terahertz frequencies, then, exists in one
of the three transport limits, in devices of half a micrometre and below
with a distant gate, with edges in the upper half of the range
measured for etched edges and with contacts that absorb electrons
arriving obliquely, and at its best it
equals the circuit term of the plasmonic element at room temperature.

\begin{table*}
\caption{Rectification by geometry and regime. $A$: ballistic asymmetry at
$0.1$~V; $D$: asymmetry of the voltage in steady flow; $\RI$: current
responsivity into a $75\ \Omega$ source; $R_{\mathrm{abs}}$: per absorbed
watt, Eq.~\eqref{eq:rabs} for the imposed fields and
Eq.~\eqref{eq:computed} for the computed one. ``Shape only'' is the symmetric termination; DS
has the drain tied to the gate.}
\label{tab:verdict}
\begin{ruledtabular}
\begin{tabular}{L{0.15\textwidth}L{0.31\textwidth}L{0.24\textwidth}L{0.23\textwidth}}
Regime & Funnel & Classical Tesla valve & What limits it \\
\hline
Ballistic, dc & $A=1.6$ to $1.9$ clean; $1.1$ to $1.6$ at $\lambda=700\nm$, by length and field profile; device D with the field computed: $1.02$ & $A=0.89$ to $1.08$, no systematic sign & reciprocity; needs the bias landscape \\
Viscous, dc, no gate & $|D-1|\leq10^{-2}$ within the drift limit & $10^{-3}$ to $3\times10^{-2}$ within the drift limit & drift speed below $0.3\,\vF$ caps $\Rey$ \\
Viscous, dc, gated & contacts exchanged: $<10^{-3}$; polarity reversed: $0.48\,V/U_{0}$ & the same polarity term; exchange $<10^{-3}$ & the polarity term is the gate, in any channel \\
Plasmonic, shape only & $10^{-4}$ to $0.13\AW$; $0.22\AW$ at taper resonance & $10^{-5}$ to $10^{-3}\AW$; $0.02\AW$ at resonance & plasmon damping; saturates first \\
Plasmonic, DS & $0.07$ to $0.34\AW$; $\eta<1\%$ & $0.02$ to $0.1\AW$ at $6$ to $11$~k$\Omega$ & $r^{2}$; gate overdrive; contacts \\
Ballistic, $1\THz$, $6\um$ & $\RI=0.05\AW$ (uniform field); $0.2$ to $0.45\AW$ (drop at neck) & none ($A=1$ at dc) & length: $fL/\vF=6$ \\
Ballistic, $1\THz$, $0.55\um$, imposed field & $R_{\mathrm{abs}}=0.7\AW$, either profile; $0.14$ to $0.42\AW$ with rough edges or $\lambda=100\nm$ & none & the field is not computed \\
Ballistic, $1\THz$, $0.55\um$, computed field & $R_{\mathrm{abs}}=0.12\AW$ with specular edges; $0.10\AW$ at $\lambda=100\nm$; rough edges: half at a head-on specularity of $0.5$, zero at $0.2$; contacts of transparency $0.25$: $0.07$ to $0.12\AW$; collimating contacts: none; $\RI=0.012$ to $0.026\AW$ & not computed & contact drops; edges; angular selection at the contacts; under a $30\nm$ gate the circuit term is larger; $\eta\sim10^{-4}$ \\
\end{tabular}
\end{ruledtabular}
\end{table*}

\begin{figure*}
\includegraphics[width=\textwidth]{\figpath/fig8_synthesis.pdf}
\caption{Synthesis. (a) Current responsivity against frequency for every
mechanism: the plasmonic $1\um$ funnel with the Dyakonov-Shur termination
and with the symmetric drive that isolates the shape, the Tesla valve, the
$6\um$ ballistic funnel with a uniform field and with the potential drop
confined to the neck, all into a $75\ \Omega$ source, and the $0.55\um$
ballistic funnel with the field computed (gate at $300\nm$, shape-only
drive), for one device on the same source and per absorbed watt. Solid
lines: $100$~K parameters or the clean ballistic limit; dashed: $300$~K
or $\lambda=700\nm$. Points consistent with zero are not drawn. (b)
Noise-equivalent power of the simulated elements, which is Johnson
limited and assumes ideal coupling, against measured graphene terahertz
detectors (Table~\ref{tab:detectors}); the ballistic array is that of
Sec.~\ref{sec:pic}.}
\label{fig:synthesis}
\end{figure*}

\subsection{What the elements are good for}

Table~\ref{tab:detectors} and Fig.~\ref{fig:synthesis}(b) place the
simulated elements against measured graphene terahertz detectors. At room
temperature the best reported devices have noise-equivalent powers of $100$
to $130\pWrtHz$ referred to the optical power at
$0.4\THz$~\cite{generalov2017graphene,ludwig2024terahertz}, with an
estimate of $51\pWrtHz$ at $0.33\THz$ for an epitaxial bilayer
device~\cite{qin2017room}, and $80$ to
$160\pWrtHz$ referred to a diffraction-limited spot at $2.5$ to
$2.8\THz$~\cite{castilla2019fast,viti2020hbn}; most others lie between
$0.5$ and
$2$~nW/Hz$^{1/2}$~\cite{spirito2014high,zak2014antenna,bandurin2018dual,tong2015antenna},
the first devices were at tens of
nW/Hz$^{1/2}$~\cite{vicarelli2012graphene}, and four-terminal ballistic
rectifiers of encapsulated graphene have been reported at
$34\pWrtHz$ below $0.7\THz$~\cite{auton2017terahertz} and, at $3\THz$ and
$4$~K, at $20$ to $30\pWrtHz$~\cite{shi2026ballistic}; funnel rectifiers of
chemical vapor deposited graphene, with a mean free path below the device
size, are at $29$~nW/Hz$^{1/2}$~\cite{hemmetter2026enhanced}. The
Johnson-limited $60\pWrtHz$ of the DS-terminated $1\um$ funnel at $300$~K,
or $150\pWrtHz$ with $300\ \Omega\um$ contacts, is therefore where the best
measured devices already are, and its voltage responsivity of $50$~V/W is
within their range of $1$ to $100$~V/W. The idealised model holds no hidden
order of magnitude; the comparison supports the model rather than
promising an improvement. The matched array of $0.55\um$ ballistic funnels
is at $130\pWrtHz$ at $1\THz$ and room temperature, twice the figure of
the plasmonic element at that temperature and again where the best
measured room-temperature devices are; it assumes nineteen devices matched
to the antenna, specular edges (at a head-on specularity of $0.5$ the
figure doubles), contacts that absorb at all angles and no parasitic
capacitance, and it is
four times the value reported for the encapsulated ballistic
rectifiers~\cite{auton2017terahertz}. At low temperature, encapsulated devices reach
$0.2\pWrtHz$ at $10$~K and
$0.13\THz$~\cite{bandurin2018resonant,gayduchenko2021tunnel} and
$18\pWrtHz$ at $77$~K and $2.8\THz$~\cite{viti2020hbn}; the simulated $17$
to $27\pWrtHz$ at $100$~K, and $1\pWrtHz$ at low density, lie in the same
range. Three caveats apply. Nearly all reported values, like ours, are
computed from the Johnson noise of the measured resistance and not from
measured noise. The reference power differs between reports (total beam,
power delivered to the antenna, or power scaled to a diffraction-limited
spot), and values normalised to the active device area are not comparable
and are left out. And part of the measured response of graphene detectors
is photothermoelectric~\cite{cai2014sensitive,castilla2019fast,ludwig2024terahertz},
which the model does not contain. The room-temperature incumbent is better
than any of these: commercial zero-bias Schottky detectors have typical
noise-equivalent powers of $5$ to $27\pWrtHz$ from $0.22$ to
$1.1\THz$~\cite{vdi2026detectors}, which the simulated plasmonic funnel
reaches only at $100$~K and the simulated ballistic array does not reach.

\begin{table*}
\caption{Measured graphene terahertz detectors and the simulated element.
Responsivity and NEP as stated by the authors; the last column gives the
power to which they are referred: the power in a diffraction-limited spot,
the full optical beam power, the power delivered to the antenna, an
extrinsic figure, the waveguide input, or, for this work, the available
power with ideal coupling or the power absorbed by a matched array. NEPs are estimated from Johnson noise except the one marked with
a double dagger, which was measured; the commercial values are typical
specifications. Entries marked with a dagger were read from the abstract
only.}
\label{tab:detectors}
\begin{ruledtabular}
\begin{tabular}{lL{0.21\textwidth}cccL{0.14\textwidth}l}
Reference & Device and mechanism & $f$ (THz) & $T$ (K) & Responsivity & NEP (pW/Hz$^{1/2}$) & Power \\
\hline
\cite{vicarelli2012graphene} & FET, overdamped plasma wave & 0.3 & 300 & 0.10 to 0.15 V/W & $3\times10^{4}$ & spot \\
\cite{spirito2014high} & bilayer FET & 0.29 to 0.36 & 300 & 1.2 V/W & $2\times10^{3}$ & spot \\
\cite{zak2014antenna}$^{\dagger}$ & CVD FET, bow-tie antenna & 0.6 & 300 & 14 V/W & 515 & optical \\
\cite{generalov2017graphene} & CVD FET, bow-tie antenna & 0.4 & 300 & 74 V/W & 130 & optical \\
\cite{ludwig2024terahertz} & CVD FET, bow-tie antenna & 0.4 & 300 & 59 to 63 V/W & 101 to 114 & optical \\
\cite{bandurin2018dual} & encapsulated FET & 0.13 & 300 & 20 V/W & 600 & antenna \\
\cite{bandurin2018resonant} & encapsulated bilayer FET & 0.13 & 10 & $>3$ kV/W & 0.2 & antenna \\
\cite{tong2015antenna} & thermoelectric, patch antenna & 1.9 & 300 & 4.9 V/W & $1.7\times10^{3}$\,$^{\ddagger}$ & optical \\
\cite{castilla2019fast} & photothermoelectric p-n junction & 2.5 & 300 & 105 V/W & 80 & spot \\
\cite{viti2020hbn} & photothermoelectric and plasma wave & 2.8 & 300; 77 & 49; 180 V/W & 160; 18 & spot \\
\cite{auton2017terahertz}$^{\dagger}$ & four-terminal ballistic rectifier & up to 0.69 & 300 & 764 V/W & 34 & extrinsic \\
\cite{qin2017room}$^{\dagger}$ & epitaxial bilayer FET & 0.33 & 300 & 30 V/W & 51 & not verified \\
\cite{shi2026ballistic} & four-terminal ballistic rectifier & 3.0 & 4; 300 & 77; 6 V/W & 20 to 30 (4 K) & spot \\
\cite{hemmetter2026enhanced} & CVD four-terminal funnel rectifier & 0.49 to 0.68 & 300 & 0.77 V/W & $2.9\times10^{4}$ & optical \\
\cite{vdi2026detectors} & commercial zero-bias Schottky & 0.22 to 1.1 & 300 & 750 to 2200 V/W & 5 to 27 & waveguide \\
\hline
This work & $1\um$ funnel, DS, simulated & 0.3 to 1 & 300 & 0.09 A/W, 52 V/W & 60; 150 with contacts & available \\
This work & same & 0.3 to 1 & 100 & 0.21 to 0.34 A/W, 36 to 58 V/W & 17 to 27 & available \\
This work & same, $5\times10^{10}\percm$ & 0.3 & 100 & & 0.9 & available \\
This work & $0.55\um$ ballistic funnel, computed field, matched array, simulated & 1 & 300 & 0.12 A/W & 130 & absorbed \\
\end{tabular}
\end{ruledtabular}
\end{table*}

For energy conversion the conclusion does not depend on such comparisons,
and its framework is not new. The efficiency of a field-effect detector is
quadratic in the amplitude up to the gate swing and has its maximum, of the
order of ten per cent, below threshold~\cite{kachorovskii2013performance};
silicon transistor rectennas work there and have reached $13.6\%$ at
$263\GHz$~\cite{khan2022cmos}. What the present calculation adds is the
graphene case. A gapless monolayer has no subthreshold regime, the voltage
scale of its gated fluid is the gate overdrive, and the efficiency at an
amplitude equal to the overdrive is the pure number $r^{2}\Rant/\Rdc$ of
Eq.~\eqref{eq:eta}. For a funnel on a real antenna that number is of order
$10^{-2}$, and no density or gate distance in the hydrodynamic window
raises it; a low density or a thin gate moves the power at which the bound
is reached from a milliwatt to a microwatt, which is useful for matching a
detector's dynamic range to its signal, and nothing more. Realistic
contacts of $100$ to $300\ \Omega\um$ cost a further factor two to five.
The ballistic funnel has a lower voltage scale, the Fermi energy, and a
lower efficiency still: $10^{-4}$ at $13\ \mu$W per device, the largest
amplitude computed. A conversion efficiency of $58\%$ at
$0.9\THz$ has been predicted by simulation for a graphene
rectenna~\cite{dragoman2016graphene}, but its rectifier is a Schottky
junction between graphene and a semiconductor, not a geometric or plasmonic
element. Within the mechanisms studied here a graphene rectenna has no
harvesting regime.

A designer who wants a graphene terahertz detector from these shapes
should therefore read the results as follows. In a gated device of a
micrometre or more, use the funnel as a matched horn, not as a diode: its
value is the impedance it presents and the taper resonance. Terminate the
channel asymmetrically, since that is where the rectification comes from,
put the saturation knee at the expected power by choice of density and
gate distance, and spend the fabrication effort on the contacts, because a
resistive drain contact removes the asymmetry it is supposed to serve. If
the device is to rectify by its shape, make it short, by
Eq.~\eqref{eq:criterion}: half a micrometre for $1\THz$, leads included.
Spend the effort on the edges before the mobility, since a mean free path
of $100\nm$ costs a fifth of the response while an rms edge roughness
of $2\nm$ costs half and $3\nm$ all of it, and establish the
specularity of the edges of the process first, since the values measured
for etched edges run from those that would leave nothing to those that
would leave four fifths; ask of the contacts that they absorb
electrons arriving at wide angles before asking for a low resistance,
since a transparency of a quarter costs little and a collimating
junction in front of the metal costs everything; keep the gate farther away than the device is long, or accept that
the gate connection does the rectifying; and provide for the impedance,
since one such funnel is a kilohm. Expect no more than the circuit term
gives. Do not judge a geometric diode for terahertz use by its dc
asymmetry, nor by a simulation with an imposed field. Do not add Tesla
loops.

\subsection{Limitations}

There is no experiment here. The solvers are validated against analytic
limits. The one measured asymmetry that can be set beside them at equal
bias, $1.03$ to $1.04$ for device D at $0.1$~V, is consistent with the
computed-field value of $1.02$, but it is too close to one, and the
measured device too much more resistive than the modelled one, to test
the model (Sec.~\ref{sec:dc}). Most idealisations favour the device. That is
the safe direction for the negative results, on the Tesla valve and on the
conversion efficiency, and the unsafe one for the positive result on the
short ballistic funnel, which computing the field reduced sixfold, which
the rougher of the etched edges that have been measured would remove, as
would contacts that collimate, and which remains an upper estimate.

The ballistic model of the $6\um$ device and of the length scaling
imposes its potential and has no displacement current. For the $6\um$
device the terahertz result depends entirely on the length of the drop,
and the scaling curve of Fig.~\ref{fig:length}(b) was confirmed with a
computed field for one device only, through its roll-off. For the
$0.55\um$ funnel the field is computed, in a model with limits of its
own. The dielectric is uniform. The contacts are equipotential
strips that do not dope the graphene beneath or beside them. Between
ideal ones a quarter of the bias drops at the contact planes within
three grid cells. The two contacts that reflect are the extremes of
angular selection, one transparency for all angles and a smooth
junction, and which of them a metal edge resembles is not known to us.
The gas loses no energy to the lattice, so that it is $3$ to $7\%$
warmer in kinetic energy between contacts that reflect at all angles and
up to $14\%$, and not yet steady, between collimating ones; the points
at a transparency of $0.75$ and those with collimating contacts are
single runs. The static pair for device D shares these limits. Its
asymmetry is the same at two particle numbers and its conductance
differs by $2.5\%$ between them; the gas was still warming at the end of
the averaging window, by one per cent of its kinetic energy across it;
and the bias of the measured $1.66$, about $1.5$~V, is beyond what a
model without optical phonons, interband processes and heat loss can be
asked. The electrons are classical particles in a classical field,
with five to seven modes in the neck, and they interact through the
common field only; short-range electron-electron scattering enters
through the $100\nm$ bound. The rough edge follows
Eq.~\eqref{eq:rough}, which is the Kirchhoff result: it ignores the
correlation length of the roughness, is least reliable towards grazing
incidence, and has not been tested against a measurement on graphene,
since none resolves the angle of incidence. Reading a focusing ratio as
the head-on specularity is our interpretation, and we found no measured
rms roughness of an etched graphene edge. All rough-edge points but one,
and the points of the single-specularity scan between $p=0.7$ and $0.9$,
are single runs. The search for measurements of edge specularity was a
bounded one and not a survey.
The antenna is a voltage source, with its
$75\ \Omega$ entering through the computed admittance; the matched array
is an idealisation, and no array or antenna was modelled. Rectified
currents below about $0.05\ \mu$A are not resolved. Two results of
Sec.~\ref{sec:pic} are reported without an explanation: the reversal of
the strip's current near the plasma resonance at $3\THz$, where three
frequencies do not resolve the resonance, and the rise of the funnel's
current above the square law at $0.2$~V. No model of the crossover between
the ballistic and the fluid limit is offered: the two terahertz models
describe the same gate and nearly the same density, $1.4\times10^{12}$
and $10^{12}\percm$, in opposite limits of the electron-electron
scattering rate, and a real device at room temperature lies between
them.

The plasmonic and the gated steady-flow models use a local gate
capacitance, which is least accurate for the $175\nm$ neck under a $30\nm$
gate; they have ideal contacts except where stated, perfect coupling to an
antenna represented by its resistance, and Johnson noise only. They contain
no electron heating, and so none of the photothermoelectric response that
contributes to measured graphene detectors; the fluid description is
quantitative to about $0.3\THz$ at $100$~K and $1\THz$ at $300$~K, where
$\omega\tauee=1$, and indicative beyond. The low-bias shape term of the
gated steady flow is bounded, not resolved. Finally, two experiments on
electronic Tesla valves are outside what the models reproduce: the
rectification reported in graphene at low
current~\cite{geurs2020rectification} is not obtained from steady
hydrodynamics with or without the gate, and the large diodicity above
threshold in GaAs~\cite{sarypov2026electron} belongs to a regime of
unsteady, heated flow that was not computed.

\section{Conclusion}
\label{sec:conclusion}

We have computed the rectification of graphene funnels and of an
electronic Tesla valve in the ballistic, viscous and plasmonic transport
limits, from dc to $30\THz$ and from the small-signal limit to the largest
amplitudes the models allow,
and separated in each what the shape does from what the circuit does.

The Tesla valve does not rectify usefully in any case computed. It shows
no ballistic dc asymmetry of a valve's sign in an imposed field, a
viscous diodicity of at most three per cent at the drift velocities the
graphene Fermi liquid allows, and a shape-only terahertz response in the
gated fluid of $0.02\AW$ or less up to $1\THz$, with a channel of
$175\nm$ or of $1\um$; it was not computed with its field or as a
ballistic device under an ac drive.

In the electron fluid the funnel's rectification belongs mainly to the
circuit. A steady current in a gated channel of any shape shows an
asymmetry of $0.48$ times the bias over the gate overdrive under reversal
of polarity and less than $10^{-3}$ under exchange of the contacts, and
the terahertz response, $0.05$ to $0.26\AW$ into a $75\ \Omega$ antenna,
is carried mainly by the Dyakonov-Shur asymmetry of the terminations;
with a viscosity that does not fall with frequency the range is $0.05$ to
$0.18\AW$. At room temperature the shape gives $2$ to $4\%$ of that
response with leads seven plasmon decay lengths long and $4$ to $54\%$
with leads of $1\um$, and a plain channel with the same terminations
gives $0.6$ to $1.3$ times the response of the funnel, according to the
leads. The response follows the square law of
field-effect detectors up to a drift speed of $0.3\,\vF$, which we take
as the limit of the model. The efficiency there is $10^{-5}$ to
$1.2\times10^{-4}$ where it was followed in time, and up to
$5\times10^{-4}$ by the square law at $100$~K if the viscosity falls with
frequency. Contact resistance attacks the asymmetry itself.

In the ballistic limit the shape rectifies, and the frequency to which it
does so is set by the length of the device: for a field spread over the
device the response is a function of $fL/\vF$, half at $0.4$ and gone at
$1.7$. The $6\um$ funnel that anchors the dc literature is then too long
by a factor of fifteen at $1\THz$. Its computed static field is
concentrated at its neck, and what it does at $1\THz$ with that field
was not computed. A $0.55\um$ funnel gives $0.7\AW$ per absorbed watt at
$1\THz$ with an imposed field and $0.12\AW$ with the field computed from
the charge of the electrons, a gate and the contacts. That response is
halved at a head-on edge specularity of $0.5$ and gone near $0.2$, inside
the range of zero to $0.8$ measured for etched
edges~\cite{taychatanapat2013electrically,masubuchi2012boundary,barnard2017absorptive,berdyugin2020minibands}.
It needs contacts that absorb electrons arriving obliquely, whatever
their resistance, and a gate farther away than a twentieth of the length
of the device.

Geometry, then, can make a graphene terahertz rectifier in two cases: the
ballistic funnel of half a micrometre on those conditions, in a model
without collisions between electrons, and the fluid funnel about one
plasmon decay length long, where the shape gives up to $0.08\AW$ at room
temperature beside terminations that give more, in the funnel or in a
strip. In neither case does it make a rectenna that harvests. What it
makes is a square-law detector with a Johnson-limited noise-equivalent
power of the order of $100$~pW/Hz$^{1/2}$ at room temperature, which is
where the best measured graphene detectors already are. Three things
would test or change these statements and are outside the models used
here: a measurement of the dc asymmetry and the terahertz response of one
short funnel against the distance of its gate, the quality of its edges
and the kind of its contacts; a rectification mechanism that is neither
ballistic, viscous nor plasmonic, of which the photothermoelectric effect
is the obvious example; and a material whose electron liquid sustains a
higher drift velocity relative to its plasma and Fermi velocities.


\begin{acknowledgments}
The simulation code, the analysis and this manuscript were developed with the
assistance of an AI coding assistant (Claude, Anthropic); the author directed
the work, reviewed the results and takes responsibility for the content. The
code, the input geometries, all result files and the scripts that produce the
figures are available in the public repository
\url{https://github.com/TtheBC01/GeometricDiodes}.
\end{acknowledgments}

\appendix
\section{Second-order plasmonic response}
\label{app:perturbation}

In units of the contact width $W$, the plasma speed $s_{0}$ and the gated
density $n_{0}$, Eqs.~\eqref{eq:cont} and \eqref{eq:mom} contain the
friction $\Gamma=W/s_{0}\taumr$, the plasma Reynolds number
$\Rey_{s}=s_{0}W/\nu$, the collision time $s_{0}\tauee/W$ and the
quantum-capacitance fraction $q$, with
$h(n)=2(1-q)(\sqrt{n}-1)+q\ln n$ and $h''(1)=-(1+q)/2\equiv\kappa$. Writing
$n=1+\mathrm{Re}(n_{1}e^{i\Omega t})+n_{20}$ and
$\bm{u}=\mathrm{Re}(\bm{u}_{1}e^{i\Omega t})+\bm{u}_{20}$, the first order is
\begin{align}
i\Omega n_{1}+\nabla\cdot\bm{u}_{1}&=0,\\
(i\Omega+\Gamma)\bm{u}_{1}+\nabla n_{1}-\nu(\Omega)\nabla^{2}\bm{u}_{1}&=0,
\end{align}
with the dimensionless viscosity
$\nu(\Omega)=\Rey_{s}^{-1}/(1+i\Omega s_{0}\tauee/W)$, and the time
average of the second order is
\begin{equation}
\nabla\cdot\bm{u}_{20}=-\tfrac12\mathrm{Re}\,\nabla\cdot(n_{1}\bm{u}_{1}^{*}),
\end{equation}
\begin{multline}
\Gamma\bm{u}_{20}+\nabla P-\Rey_{s}^{-1}\nabla^{2}\bm{u}_{20}\\
=-\tfrac12\mathrm{Re}\,(\bm{u}_{1}^{*}\cdot\nabla)\bm{u}_{1}
+\tfrac14\mathrm{Re}\,\bm{u}_{1}^{*}(\nabla\cdot\bm{u}_{1})
+\tfrac18\nabla|\bm{u}_{1}|^{2},
\end{multline}
where $P=n_{20}+(\kappa/4)|n_{1}|^{2}$; the last two terms on the right
are absent for the classical gas. The dc contact voltage is
$P+c_{V}|n_{1}|^{2}$ with $c_{V}=1/8$ for the Dirac gas and zero for
the classical gas, and the dc particle current through a contact is the flux of
$\bm{u}_{20}+\mathrm{Re}(n_{1}\bm{u}_{1}^{*})/2$. Each contact is an
equipotential with one unknown voltage per order, closed by Kirchhoff's law
for its termination: $R I_{c}+V_{c}=V_{a,c}$ for a source of resistance $R$,
$I_{c}=0$ for an open contact. The same structure with a matched load at dc
gives the rectified power. Voltages are in units of
$V_{u}=m^{*}s_{0}^{2}/e$, resistances in units of
$m^{*}s_{0}/e^{2}n_{0}W$ ($65\ \Omega$ for $W=4\um$), and frequency in units
of $s_{0}/2\pi W$ ($0.142\THz$).

For a uniform channel of length $L$ with an ideal source of amplitude
$U_{a}$ and an open drain these equations integrate to
\begin{multline}
\frac{\Delta U}{U_{a}^{2}}=\frac{\Phi''-\kappa}{4}\left(|n_{1}(L)|^{2}-1\right)
+\frac{\alpha}{4}|u_{1}(0)|^{2}\\
+\frac{\Gamma}{2}\int_{0}^{L}\mathrm{Re}\,(n_{1}u_{1}^{*})\,dx ,
\label{eq:channel}
\end{multline}
with the first-order fields per unit source amplitude,
$n_{1}=\cos k(L-x)/\cos kL$ and $u_{1}=-n_{1}'/(i\Omega+\Gamma)$,
$k^{2}=\Omega^{2}-i\Omega\Gamma$, $\Phi''(1)=-q/2$, and $\alpha$ the
coefficient of $u\,\partial_{x}u$ in one dimension: $\alpha=0$ for
Eq.~\eqref{eq:mom} and $\alpha=1$ for the ordinary convective term. For
the classical gas, $\alpha=1$ and
$\kappa=\Phi''=0$, this is the result of Dyakonov and
Shur~\cite{dyakonov1996detection}. With $\alpha=1/2$ and the density fixed
at the source, Eq.~\eqref{eq:channel} is Eq.~(43) of
Ref.~\cite{tomadin2013theory}, which we reproduce to $3\times10^{-10}$.
The lines of Fig.~\ref{fig:small}(a) are Eq.~\eqref{eq:channel} for the
classical gas and for the Dirac gas, $\alpha=0$.

The integral can be removed. The first-order equations give
$\partial_{x}(|n_{1}|^{2}+|u_{1}|^{2})=-2\Gamma\,\mathrm{Re}\,(n_{1}u_{1}^{*})$,
and with $n_{1}(0)=1$ and $u_{1}(L)=0$
\begin{equation}
\frac{\Delta U}{U_{a}^{2}}=\frac{1+\alpha}{4}|u_{1}(0)|^{2}
-\frac{1-\Phi''+\kappa}{4}\left(|n_{1}(L)|^{2}-1\right).
\label{eq:half}
\end{equation}
For the classical gas the two coefficients are $\tfrac12$ and $\tfrac14$;
for the Dirac gas $1-\Phi''+\kappa=\tfrac12$ for any $q$, and they are
$\tfrac14$ and $\tfrac18$. In units of its own $V_{u}$ the Dirac gas
therefore gives exactly half the response of the classical gas, at every
frequency, friction and length. The statement extends to two dimensions
when there is no viscosity: the first-order velocity is then a gradient,
and the time-averaged equations of the Dirac gas are those of the
classical gas with the dc current and the dc voltage halved, for any
shape and any linear terminations. With viscosity it does not hold, and
the funnels of Sec.~\ref{sec:small} are not at half the classical
response.

Beyond second order the momentum per particle of the drifting gas is
$m^{*}(n)\gamma^{1/2}\bm{u}$ and its pressure is
$P_{0}(n)\gamma^{-3/2}$, with $\gamma=(1-u^{2}/\vF^{2})^{-1/2}$ and $n$ the
density in the frame of the lattice, and the momentum balance becomes
\begin{multline}
\Bigl(1+\frac{\gamma^{2}}{2\vF^{2}}\bm{u}\bm{u}^{T}\Bigr)
\bigl[\partial_{t}\bm{u}+(\bm{u}\cdot\nabla)\bm{u}\bigr]
-\tfrac12\bm{u}\,(\nabla\cdot\bm{u})-\tfrac14\nabla u^{2}
+\frac{\bm{u}}{\taumr}\\
=-\frac{1}{m^{*}}\Bigl[\gamma^{-1/2}\,\nabla\frac{e^{2}n}{C_{g}}
+\gamma^{-2}\,\nabla\bigl(\hbar\vF\sqrt{\pi n}\bigr)\Bigr]
+\nu\nabla^{2}\bm{u}.
\label{eq:exact}
\end{multline}
This is the Euler equation of an ideal relativistic fluid with $\vF$ for
the speed of light~\cite{landau1987fluid}, multiplied by the matrix
$1+\gamma^{2}\bm{u}\bm{u}^{T}/\vF^{2}$, with the gate force and the
friction added; in the units used here $u^{2}/\vF^{2}=u^{2}/2q$, so no new
parameter enters. It reduces to Eq.~\eqref{eq:mom} at small $u$. On an
ungated fluid that drifts at $0.05$ to $0.8\,\vF$ its wave speeds obey
the relativistic addition of velocities to six digits, and a steady
radial flow computed with it agrees with the radial momentum balance
written directly from the stress tensor to $3\times10^{-4}$ of the
inertial part of the voltage. The friction and the viscous term are kept
in the form of Eq.~\eqref{eq:mom}; the contact condition is the potential
$\Phi(n)$ at the density $n$.

A contact resistance $R_{c}$ enters as a change of the ports only: the
source becomes a Thevenin source of resistance $\Rant+R_{c}$ and the drain
one of zero amplitude and resistance $R_{c}$. In second order the dc
problem is linear, so the short-circuit current and open-circuit voltage
satisfy $I_{\mathrm{sc}}(\Rdc+2R_{c})=V_{\mathrm{oc}}$ exactly, which is
one of the ladder tests. At the signal frequency the antenna does not see
the device impedance plus $2R_{c}$, because the source and drain currents
differ by the gate displacement current.

\section{The drift limit on the density and gate grid}
\label{app:knee}

The powers and efficiencies of Figs.~\ref{fig:operating}(a) and (b) come
from the small-signal solver alone. At each grid point the largest
first-order drift speed in the device, $\max|\bm{u}_{1}|$ per unit antenna
amplitude, gives the amplitude $V_{\mathrm{lim}}$ at which the drift
reaches $0.3\,\vF$, the available power
$P_{\mathrm{lim}}=V_{\mathrm{lim}}^{2}/8\Rant$, and the efficiency
$(\eta/\Pav)\,P_{\mathrm{lim}}$ of the square law there. The time-domain
runs support the two extrapolations in amplitude: for the reference
device the largest drift speed
is linear in the amplitude to $9\%$ up to $0.3\,\vF$ and the efficiency
follows the square law of the same solver to $7\%$, and a second ladder
at $(5\times10^{11}\percm,10\nm)$ gives $12\%$ and $9\%$. They do not
confirm the level at $100$~K. There the time-domain solver, which has no
viscoelastic term, finds a drift per unit amplitude smaller by factors of
$1.1$, $1.5$ and $2.5$ at $0.1$, $0.3$ and $1\THz$ and a different square
law (Sec.~\ref{sec:large}); with the viscosity held constant in the
small-signal solver the two agree on the drift to $1\%$ at $0.1$ and
$0.3\THz$ and $5\%$ at $1\THz$, and on the law to the error of the
coarser mesh, $6$ to $16\%$. The grid values are
those of the small-signal solver with the viscoelastic term. At $300$~K the two agree on the drift to $3\%$ and on the law to
$1\%$ at $0.1$ and $0.3\THz$.

\section{Steady gated flow in one dimension}
\label{app:gated}

In a gradual channel of width $W(x)$ carrying the current $I$ the
continuity equation gives $n u W=I$. Without the convective and viscous
terms the steady form of Eq.~\eqref{eq:mom} is $\Gamma u=-h'(n)\,dn/dx$ in
the units of Appendix~\ref{app:perturbation}, so that
\begin{equation}
n\,h'(n)\,\frac{dn}{dx}=-\frac{\Gamma I}{W(x)} .
\end{equation}
With $G'(n)=n h'(n)$, i.e. $G=n^{2}/2$ for the classical gas and
$G=\tfrac23(1-q)n^{3/2}+qn$ for the Dirac gas, integration from the
driven contact, where $n=n_{s}$, to the contact tied to the gate, where
$n=1$, gives Eq.~\eqref{eq:1d}. The voltage of the driven contact is
$\Phi(n_{s})$ with $\Phi=(1-q)(n-1)+2q(\sqrt{n}-1)$ in units of $V_{u}$.
The width enters only through $\int dx/W$, which is unchanged when the
channel is turned end for end, so exchanging the contacts cannot produce
an asymmetry. Reversing the current changes the sign of the right-hand
side; expanding to second order with $G''(1)=(1-q)/2$ and
$\Phi''(1)=-q/2$,
\begin{equation}
V=\rho I-\tfrac14(\rho I)^{2},\qquad \rho=\Gamma\int_{0}^{L}\frac{dx}{W},
\end{equation}
and the ratio of the two voltage magnitudes is $1+\rho|I|/2=1+V/2V_{u}$,
which is Eq.~\eqref{eq:dpol} with $U_{0}=(1-q)V_{u}$ and $q=0.039$; for
the classical gas the coefficient of $(\rho I)^{2}$ is $\tfrac12$ and the ratio
is $1+V/U_{0}$. In a uniform channel the inertia can be kept. For the
classical gas the equation becomes $(nh'-j^{2}/n^{2})\,dn/dx=-\Gamma j$
with $j=I/W$, with the integral $G(n)+j^{2}/n$. For the Dirac gas the
three inertial terms of Eq.~\eqref{eq:mom} cancel in one dimension and
what remains of Eq.~\eqref{eq:exact} is of relative order $u^{2}/\vF^{2}$;
it is integrated numerically. The two-dimensional solver reproduces both
to $1.1\times10^{-4}$ for both signs up to $V=0.3\,U_{0}$, and the Dirac
gas to $7\times10^{-5}$ up to a drift of $0.8\,\vF$. In a channel of
varying width the inertia brings in $dW/dx$, and with it the inertial
shape term of Sec.~\ref{sec:dc}, which is second order in the current.



\bibliography{refs,refs_detectors,refs_prior}

\end{document}
